\documentclass[10pt,a4paper]{article}
\usepackage[numbers,sort&compress]{natbib}
\usepackage[english,french]{babel}
\usepackage{popsnotes}
\selectcolormodel{natural}
\usepackage{algpseudocode}
\algrenewcommand\algorithmicwhile{\textbf{tant que}}
\algrenewcommand\algorithmicfor{\textbf{pour}}
\algrenewcommand\algorithmicforall{\textbf{pour chaque}}
\algrenewcommand\algorithmicdo{\textbf{faire}}
\algrenewcommand\algorithmicend{\textbf{fin}}
\algrenewcommand\algorithmicreturn{\textbf{renvoyer}}
\algrenewcommand\algorithmicloop{\textbf{boucle}}
\algrenewcommand\algorithmicif{\textbf{si}}
\algrenewcommand\algorithmicthen{\textbf{alors}}
\geometry{left=23mm,right=23mm,top=20mm,bottom=20mm,headsep=4mm}
\setlength{\headwidth}{\textwidth}
\setcounter{secnumdepth}{-1}
\setcounter{tocdepth}{1}
\addto\captionsfrench{\renewcommand{\contentsname}{Sommaire}\renewcommand{\figurename}{Figure}}
\titleformat{\section}{\normalfont\rmfamily\bfseries\Large\color{black}}{}{0pt}{}[\vspace{2pt}\titlerule]
\titleformat{\subsection}{\normalfont\rmfamily\bfseries\large\color{black}}{}{0pt}{}[\titlerule]
\newcommand{\DocHeader}{Distance d'édition sur ROMEO}
\newcommand{\reposource}[1]{\href{https://github.com/wolf75222/nvpd-edit-distance/blob/main/#1}{\nolinkurl{#1}}}
\newcommand{\sourcefig}[1]{\par\vspace{1pt}{\raggedright\footnotesize Code\,: \reposource{#1}\par}}
\newcommand{\chart}[4]{\begin{figure}[H]\centering
  \includegraphics[width=.93\linewidth,height=.35\textheight,keepaspectratio]{../figures/#1}
  \caption{#2}\label{#3}\sourcefig{#4}\end{figure}}
\newcommand{\schema}[4]{\begin{figure}[H]\centering
  \input{figures/#1}\caption{#2}\label{#3}\sourcefig{#4}\end{figure}}
\newcommand{\synthese}[2]{\par\medskip\noindent\textbf{\underline{Synthèse.}}\ #1\newline #2\par}
\newcommand{\methode}[2]{\par\smallskip\noindent\textbf{Algorithme #1.} \emph{#2}\par}
\newcommand{\chapitre}[1]{\clearpage\section{#1}}
\captionsetup{font=small,labelfont=bf,labelsep=period,skip=3pt}
\setlength{\intextsep}{6pt plus 1pt minus 1pt}
\hypersetup{pdftitle={Parallel Edit Distance for Genomic Sequences},pdfsubject={Alignement global sur ROMEO}}

\begin{document}
\thispagestyle{firstpage}
\begin{center}
{\Large\bfseries Parallel Edit Distance for Genomic Sequences\par}
\vspace{4pt}\hrule\vspace{8pt}
Master CHPS \quad Université de Reims Champagne-Ardenne\par
CHPS1001 \quad 24 novembre 2026
\end{center}
\noindent\textbf{Résumé.} La distance d'édition globale compte les opérations qui transforment une séquence entière en une autre. NvPD répartit les colonnes d'une même ligne entre cœurs ou fils GPU. Pour une collection, les paires indépendantes offrent un second niveau de parallélisme. Les mesures sur EPYC et GH200 relient le temps de calcul aux synchronisations et à la divergence des séquences. Les applications portent sur des assemblages humains et des arbres de distance.
{\small\setlength{\parskip}{0pt}\tableofcontents}
\section{Du problème biologique aux noyaux de calcul}
Une substitution remplace une base. Une insertion ou une délétion décale toutes les positions qui la suivent. La distance de Hamming compte les différences position par position pour deux chaînes de même longueur. La distance d'édition autorise ces décalages. Elle vérifie
\[
 |m-n|\le d(A,B)\le\max(m,n),\qquad
 d(A,C)\le d(A,B)+d(B,C).
\]
La première borne vient du changement de longueur. Pour la seconde borne sur \(d(A,B)\), il suffit de substituer au plus \(\min(m,n)\) bases et d'ajuster la longueur. L'inégalité triangulaire résulte de la composition de deux suites d'éditions. Chaque édition admet une opération inverse de même coût, d'où la symétrie. Plusieurs alignements peuvent réaliser le même minimum.

Une édition compte pour une unité, y compris dans les répétitions. Un écart de dix bases coûte donc dix. Cette convention répond à une question algorithmique précise. Pour interpréter des différences biologiques, il faut aussi tenir compte de l'orientation des brins, des ambiguïtés et des remaniements. Un nombre d'éditions entre deux assemblages inclut leurs différences de représentation. Il ne donne à lui seul ni un nombre de mutations historiques ni un temps de divergence.

\subsection{Calcul dense, coût et alignement}
La programmation dynamique (DP) classique utilise
\[
 C_{ij}=\min(C_{i-1,j}+1,\ C_{i,j-1}+1,\ C_{i-1,j-1}+[A_i\ne B_j]).
\]
Le parcours séquentiel avance par lignes. Le parcours parallèle avance par anti-diagonales de somme \(s=i+j\), avec \(\max(1,s-n)\le i\le\min(m,s-1)\). Chaque case lit deux valeurs de l'anti-diagonale \(s-1\) et une de \(s-2\). Les anti-diagonales d'extrémité n'occupent qu'une fraction des fils.

\begin{center}\begin{minipage}{.94\linewidth}
\textbf{Algorithme 1.} \textit{NvPD pour le coût global.} Entrées \(A,B\). Sortie \(d(A,B)\).
\begin{algorithmic}[1]
\ForAll{\(c\in\mathrm{alphabet}(A)\)}
  \State \(K[c,0]\gets 0\) ; pour \(j=1,\ldots,n\), \(K[c,j]\gets j\) si \(B_j=c\), sinon \(K[c,j-1]\)
\EndFor
\State \(P[0:n]\gets (0,1,\ldots,n)\)
\For{\(i=1,\ldots,m\)}
  \State \(Q[0]\gets i\)
  \State \textbf{en parallèle pour} \(j=1,\ldots,n\)
  \State \hspace{1em}\(k\gets K[A_i,j]\)
  \State \hspace{1em}\(Q[j]\gets\min\{1+\min(P[j-1],P[j]),\ P[\max(0,k-1)]+j-k\}\)
  \State Synchroniser les colonnes et échanger \(P,Q\)
\EndFor
\State \Return \(P[n]\)
\end{algorithmic}
\end{minipage}\end{center}
Le temps de calcul dense est \(\Theta(mn)\), auquel s'ajoute la construction de l'index en \(O(qn+m)\). Pour le coût seul, placer la chaîne courte sur les lignes réduit le nombre de synchronisations. En mode trace, le programme conserve l'orientation des entrées et remonte dans la matrice complète, selon la récurrence classique. Un test rejoue chaque opération du CIGAR\footnote{Suite de longueurs et d'opérations. Les symboles \texttt{=}, \texttt{X}, \texttt{I} et \texttt{D} représentent une égalité, une substitution, une insertion et une délétion.}, vérifie la consommation des deux chaînes et recompte le coût.

\newpage
\section{Architectures des versions}
\begin{figure}[H]\centering
\begin{tikzpicture}[font=\small\rmfamily,
  block/.style={draw=black!65,rounded corners=1pt,align=center,
    inner xsep=4pt,inner ysep=3pt,minimum height=9mm},
  flow/.style={-{Latex[length=1.5mm]},draw=black!70,line width=.5pt}]
\node[block,fill=powderGrey!18] (fasta) at (0,0) {FASTA};
\node[block,fill=powderBlue!25] (local) at (2.9,.77) {\texttt{editdist}};
\node[block,fill=powderBlue!25] (mpi) at (2.9,-.77) {\texttt{editdist-mpi}};
\node[block,fill=powderGold!20,text width=5.2cm] (kernels) at (7.25,.77)
  {DP, NvPD, edlib et fronts d'onde\\Calcul sur processeur ou accélérateur};
\node[block,fill=powderGold!20,text width=5.2cm] (ranks) at (7.25,-.77)
  {NvPD sur plusieurs rangs MPI\\Répartition des colonnes ou des paires};
\node[block,fill=powderTeal!24,text width=3.0cm] (score) at (12.1,.77)
  {Score, CIGAR\\ou matrice};
\node[block,fill=powderTeal!24,text width=3.0cm] (scores) at (12.1,-.77)
  {Score, liste de scores\\ou matrice};
\draw[draw=black!70] (fasta.east)--(1.18,0) (1.18,.77)--(1.18,-.77);
\foreach \id in {local,mpi}{\draw[flow] (1.18,0 |- \id.west)--(\id.west);}
\draw[flow] (local)--(kernels); \draw[flow] (mpi)--(ranks);
\draw[flow] (kernels)--(score); \draw[flow] (ranks)--(scores);
\end{tikzpicture}

\caption{Chemins d'exécution. Le lecteur valide les identifiants et l'alphabet. Les exécutables échangent leurs résultats par fichiers pour vérifier les distances avec des implémentations indépendantes.}
\end{figure}
\texttt{src/fasta.cpp} lit les enregistrements une fois. Le mode \texttt{pair} retourne le coût, le temps et le pic RSS\footnote{Le RSS maximal est la plus grande quantité de mémoire physique résidente du processus. Il inclut le programme, les bibliothèques, leurs espaces de travail et les entrées.} en JSON. Le mode \texttt{matrix} parcourt le triangle supérieur, recopie sa partie symétrique et écrit un TSV étiqueté. Pour les méthodes séquentielles, OpenMP peut distribuer les paires. Pour NvPD, les fils coopèrent déjà au sein d'une paire. \texttt{dp.cpp} partage l'index et la reconstruction avec CUDA et Kokkos. Les appels à edlib et à WFA2 passent par des adaptateurs. QuickEd dispose d'un exécutable séparé pour isoler les symboles C de ses dépendances.

\subsection{Réduire le domaine exploré}
Un chemin de coût au plus \(k\) vérifie \(\lvert i-j\rvert\le k\) en toute case. Le noyau \texttt{banded} applique la DP dans cette bande, selon le principe du calcul à seuil d'Ukkonen~\cite{ukkonen}. Il répète le calcul avec un seuil doublé jusqu'à obtenir le coût exact. Si le minimum d'une ligne excède le seuil, aucun chemin admissible ne peut la franchir. La version \texttt{banded-trim} retire les préfixes et suffixes communs. Cette simplification découle de \(d(cx,cy)=d(x,y)\) et de la même identité sur les chaînes renversées.
\begin{center}\begin{minipage}{.94\linewidth}
\textbf{Algorithme 2.} \textit{Distance exacte à bande adaptative.}
\begin{algorithmic}[1]
\State Si demandé, retirer les extrémités communes ; traiter les chaînes vides
\State \(k\gets\max(1,|m-n|)\)
\Loop
  \State Initialiser les bords de la bande et ses sentinelles à l'extérieur
  \State Calculer \(C_{ij}\) pour \(\max(0,i-k)\le j\le\min(n,i+k)\) dans deux lignes
  \If{la dernière case est atteinte et \(C_{mn}\le k\)} \State \Return \(C_{mn}\) \EndIf
  \State \(k\gets\min(2k,\max(m,n))\)
\EndLoop
\end{algorithmic}
\end{minipage}\end{center}
Les largeurs successives forment une suite géométrique. Avec \(L=\max(m,n)\) et \(d=d(A,B)\), la borne en temps est \(O(L(d+1))\), la mémoire de travail \(O(d+1)\). Le test de seuil emploie des entiers. edlib regroupe les différences de plusieurs cases dans les bits d'un mot machine~\cite{myers,edlib}. WFA2-lib conserve, pour chaque coût et chaque diagonale atteinte, la position la plus avancée. Chaque front prolonge gratuitement les égalités après une édition~\cite{wfa}. L'adaptateur choisit la métrique d'édition globale, en mode coût seul, sans heuristique. QuickEd estime d'abord une borne supérieure \(\hat s\), puis résout exactement dans le domaine ainsi restreint~\cite{quicked}. Son chemin par défaut reconstruit aussi l'alignement.

\chapitre{Effet de la divergence à grande longueur}
Le générateur applique des éditions à une chaîne de référence et conserve la graine de chaque paire. Le taux injecté compte les opérations du générateur. La distance \(d\) minimise le coût parmi tous les chemins et peut être inférieure à ce nombre. Au score \(s\), WFA conserve la plus grande colonne \(F_s(k)\) atteinte sur la diagonale \(k=j-i\). Prolonger une correspondance ne consomme aucun point de score~\cite{wfa}.
\methode{3}{Front d'édition à coûts unitaires.}
\begin{algorithmic}[1]
\State Étendre le préfixe commun sur \(k=0\) au score nul
\While{le point \((m,n)\) reste hors du front}
\State Augmenter \(s\) de un et examiner les diagonales admissibles \(|k|\le s\)
\State Pour chaque diagonale, retenir le plus avancé des prédécesseurs admissibles pour une insertion, une délétion ou une substitution
\State Étendre les lettres égales depuis ce point et mémoriser la colonne terminale
\EndWhile
\State \Return \(s\)
\end{algorithmic}
Pour les prédécesseurs intérieurs, les trois colonnes candidates valent \(F_{s-1}(k-1)+1\), \(F_{s-1}(k+1)\) et \(F_{s-1}(k)+1\). Chaque mouvement reste dans le rectangle des préfixes. Au plus \(2s+1\) diagonales sont actives au score \(s\), soit \(O((d+1)^2)\) entrées de front jusqu'au score final, auxquelles s'ajoutent les comparaisons des lettres égales.
\chart{01_temps_distance.png}{Les temps sur les paires synthétiques de la campagne initiale sont rapportés à la longueur et au taux d'éditions injectées.}{fig:temps}{bench/run.py}
À longueur fixée, NvPD parcourt toujours \(mn\) cases. La figure~\ref{fig:temps} montre que le temps de WFA2 augmente avec la divergence. Le débit dense \(mn/(10^9 T)\) compte des milliards de cases utiles par seconde pour NvPD, hors recalcul des halos. Pour un front, le même quotient est un débit équivalent.
Sur les cinq paires de chaque taux à un million de bases, le temps médian de WFA2 passe de 0,61 à 11,64\,s entre 2 et 10\,\% d'éditions injectées. Sur les mêmes groupes, edlib passe de 1,71 à 6,03\,s.
\synthese{Multiplier par cinq le taux injecté multiplie ici le temps médian de WFA2 par environ dix-neuf.}{La comparaison des deux décompositions denses conserve le même nombre de cases utiles par entrée.}
\chapitre{Comparaison des décompositions OpenMP à grande largeur}
Le nombre de synchronisations et la largeur des étapes distinguent les deux décompositions. Pour une matrice de \(m\times n\) cases, les anti-diagonales ont \(m+n-1\) niveaux intérieurs de largeur variable. NvPD traite \(m\) lignes de largeur \(n\) et lit un prédécesseur indirect dans l'index des occurrences.

La comparaison OpenMP reprend les quinze entrées d'un million de bases avec l'exécutable de la campagne initiale, antérieure à la vectorisation. Chaque version emploie huit cœurs. À chaque vague, trois processus par nœud traitent les trois taux sur des cœurs distincts. Les nœuds échangent les algorithmes entre deux vagues. Les cinq valeurs d'un taux correspondent à cinq paires différentes.
\chart{13_openmp_un_million.png}{Les durées OpenMP sur quinze entrées et le gain SIMD sur trois de ces entrées proviennent de deux campagnes distinctes à huit cœurs par processus.}{fig:million}{bench/complete_1m.sbatch}
Sur le panneau a) de la figure~\ref{fig:million}, les médianes des anti-diagonales vont de 190,04 à 190,29\,s selon le taux. Les médianes de NvPD vont de 309,96 à 310,91\,s. Les trente scores concordent avec les deux bibliothèques.
\synthese{À 5\,\% d'éditions injectées, la version initiale de NvPD prend 1,63 fois le temps des anti-diagonales.}{Le panneau b) isole l'effet de la vectorisation de sa boucle de colonnes.}


\chapitre{Vectorisation des colonnes sur processeur}
La boucle interne de NvPD calcule une ligne à partir de la table des dernières occurrences. Le branchement sur \(k=0\), \(k=j\) ou \(0<k<j\) peut être remplacé par le minimum de l'algorithme 1. Quand \(k=j\), le second terme vaut \(P_{j-1}\), le coût d'une correspondance. Le cas sans correspondance suit la borne établie pour l'algorithme 1. Une instruction SIMD peut alors traiter plusieurs colonnes qui choisissent des prédécesseurs différents.
\chart{03_variantes_cpu.png}{Les rapports de temps comparent les noyaux sur des paires identiques à largeur de ligne et nombre de cœurs fixés.}{fig:simd}{bench/profile_cpu_driver.cpp}
Les essais de la figure~\ref{fig:simd} utilisent GCC~11.5 avec \texttt{-O3 -march=native} sur EPYC~9654. La médiane des rapports \(T_\mathrm{initial}/T_\mathrm{SIMD}\) porte sur trois répétitions appariées. Le compilateur vectorise la boucle sur 32 octets, soit huit entiers par instruction vectorielle. À un cœur, le gain atteint 5,85 sur la forme \(256\times10^6\). Dans le panneau b) de la figure~\ref{fig:million}, la boucle initiale et la boucle SIMD utilisent les mêmes huit cœurs et les mêmes options de compilation. À 2\,\%, le temps passe de 324,41 à 53,04\,s pour \(10^6\times1\,000\,033\) cases utiles, soit 18,85\,GCUPS pour le noyau SIMD. Pour les trois taux, les gains vont de 6,12 à 6,29 avec un seul passage de chaque variante par paire.
\chart{05_cout_instructions.png}{Les compteurs de l'EPYC rapportent les instructions retirées et les erreurs de prédiction au nombre de cases des trois matrices.}{fig:instructions}{bench/profile_cpu.py}
La figure~\ref{fig:instructions} rapporte les compteurs au travail dense \(W=mn\). Sur un seul cœur et \(20\,000^2\) cases, la reformulation ramène les instructions retirées de 21,76 à 2,13 par case. Les cycles comprennent aussi les dépendances de données et les attentes mémoire.
\synthese{Sur le carré de 20\,000 bases, la reformulation accélère le calcul à un cœur d'un facteur 2,32.}{La lecture indirecte en \(k-1\) conduit à distinguer l'effet des instructions de celui du cache.}

\chapitre{Empreinte de l’index et localité sur processeur}
La taille de l'index détermine la part du jeu de données qui tient en cache. Pour \(q\) lettres distinctes et \(n\) colonnes, l'index et les deux lignes du noyau de coût seul occupent
\[
 M_{\mathrm{NvPD}}=4(q+2)(n+1)\quad [\text{octets}].
\]
Les deux lignes sont contiguës. Le dernier prédécesseur varie avec la lettre. Le prototype compact code l'écart dans un octet lorsque sa valeur et les sentinelles tiennent dans cette représentation. La conversion garde provisoirement l'index entier. Chaque nœud CPU réunit deux EPYC 9654 de 96 cœurs. Huit cœurs partagent 32\,Mio de cache L3. Chaque cœur dispose de 32\,Kio de L1 de données et de 1\,Mio de L2.
\chart{04_memoire_cpu.png}{Le pic de mémoire et les défauts L1 de lecture sont comparés aux mêmes dimensions de matrice sur huit cœurs.}{fig:mem}{bench/profile_cpu.py}
Sur la forme \(64\times4\,194\,304\), la figure~\ref{fig:mem} distingue le maximum d'espace de travail, qui passe de 96 à 80\,Mio, du RSS\footnote{Le maximum de mémoire physique résidente comprend le programme, ses entrées et les allocations des bibliothèques.}. Sur le carré à un cœur, le nombre de défauts L1 pour mille cases varie de 188,97 à 189,39 entre la boucle initiale et la boucle SIMD. Un défaut L1 peut être servi par le L2 ou le L3.
\synthese{Sur la forme la plus large, l'index compact ramène le RSS d'environ 101,5 à 85,0\,Mio.}{Le rapport des temps SIMD/compact va de 0,88 à 1,14 sur les neuf configurations, ce qui motive l'étude du placement des fils.}

\chapitre{Passage à l’échelle en mémoire partagée}
Le passage à l'échelle forte répartit un travail constant entre \(p\) cœurs. Le passage à l'échelle faible garde le nombre de cases par cœur constant~\cite{amdahl,gustafson}. Pour une matrice de travail \(W\), les indicateurs valent
\[
 S_p=\frac{T_1(W)}{T_p(W)},\qquad E_p=\frac{S_p}{p},\qquad
 E_p^{\mathrm{faible}}=\frac{T_1(W_1)}{T_p(pW_1)}.
\]
La série forte utilise \(20\,000^2\) cases. La série faible conserve 4\,096 lignes pour \(4\,096p\) colonnes. Elle fixe donc le travail à 16\,777\,216 cases par cœur et élargit les tableaux avec \(p\).
\chart{02_echelle_cpu.png}{Les courbes donnent l'accélération forte et l'efficacité faible de NvPD selon le nombre de cœurs physiques de l'EPYC.}{fig:echellecpu}{bench/profile_cpu.sbatch}
À taille fixe, huit cœurs prennent 27,93\,ms contre 31,28\,ms à seize sur la figure~\ref{fig:echellecpu}. Au dernier placement, les cœurs occupent deux groupes de cache et deux domaines NUMA.\footnote{La distance entre les cœurs et les contrôleurs mémoire dépend du domaine NUMA auquel chacun appartient.} Les synchronisations entre lignes et les accès mémoire interviennent dans le même temps mesuré.
\synthese{Huit cœurs atteignent 73,6\,\% d'efficacité forte sur la paire carrée.}{L'efficacité faible descend à 62,4\,\% à seize cœurs, avec des lignes seize fois plus larges qu'à un cœur.}

\section{Deux organisations du front de calcul}
La version CUDA partage le temps de calcul entre les lancements et les accès à la mémoire globale. Soit \(J=[u,v]\) le bloc de colonnes de sortie et \(h\le 8\) le nombre de lignes de la tuile. À la sous-étape \(t\), le domaine utile est \([\max(0,u-(h-t)R),v]\). Cette définition propage exactement les dépendances, par récurrence sur la profondeur. Les valeurs à gauche de la première colonne sont traitées par les conditions aux bords.
\begin{center}\begin{minipage}{.94\linewidth}
\textbf{Algorithme 3.} \textit{Pas CUDA de \(h\) lignes sur une bande de sortie.}
\begin{algorithmic}[1]
\State Calculer \(R\) à partir des dernières occurrences et choisir la largeur des blocs
\For{chaque groupe de \(h\) lignes successives}
\State \textbf{en parallèle pour chaque bloc} \(J\)
\State \hspace{1em}Lire les prédécesseurs de la première sous-étape dans la ligne globale
\State \hspace{1em}\textbf{pour} \(t=1,\ldots,h\), calculer les colonnes utiles par NvPD
\State \hspace{2em}Conserver les résultats intermédiaires en mémoire partagée
\State \hspace{2em}Synchroniser les fils du bloc avant la sous-étape suivante
\State \hspace{1em}Écrire \(J\) sur la nouvelle ligne globale, sans écrire son halo
\EndFor
\end{algorithmic}
\end{minipage}\end{center}
\subsection{Réutilisation des tampons entre paires}
Sur un lot sériel de paires, les deux tableaux GPU peuvent conserver leur allocation jusqu'à la paire suivante. Le programme réalloue chaque tableau uniquement si sa taille change. Les essais instrumentent les appels \texttt{cudaMalloc} et \texttt{cudaFree} du même noyau. Ils incluent le premier remplissage des tampons et leur libération finale. Chaque score est comparé à edlib. Sur les 2\,016 paires de phage, le nombre d'allocations physiques passe de 4\,032 à 250. Le pic réservé reste de 150\,892 octets. Les temps médians passent de 11,804 à 8,549\,s, avec une accélération appariée de 1,38. Sur les lots synthétiques de 256, 5\,000 et 20\,000 bases, les accélérations appariées sont respectivement de 1,59, 1,05 et 1,45. Le mode matrice réutilise ces deux tampons à l'échelle du lot.

\subsection{Fronts par coût d'édition}
Pour le calcul par coût, le front est une collection de diagonales. Sur \(k=i-j\), on note \(F_s(k)\) la plus grande position \(i\) atteinte au coût \(s\) après prolongement des égalités. Les trois choix insertion, substitution et délétion donnent respectivement \(F_{s-1}(k+1)\), \(F_{s-1}(k)+1\) et \(F_{s-1}(k-1)+1\). Seules les transitions issues d'une position valide et réalisables dans les deux chaînes sont admissibles. Cette écriture correspond au cas de coût unitaire de WFA~\cite{wfa}.
\begin{center}\begin{minipage}{.94\linewidth}
\textbf{Algorithme 4.} \textit{Front d'onde exact, coût unitaire.}
\begin{algorithmic}[1]
\State \(F_0(0)\gets\) longueur du préfixe commun, \(s\gets0\)
\While{le front ne contient pas \((m,n)\)}
\State \(s\gets s+1\)
\For{chaque diagonale atteignable \(k\) avec \(|k|\le s\)}
\State Prendre la plus grande position admissible issue du front précédent
\State Prolonger sur la diagonale tant que les deux bases concordent
\State Enregistrer la nouvelle position dans \(F_s(k)\)
\EndFor
\EndWhile
\State \Return \(s\)
\end{algorithmic}
\end{minipage}\end{center}
Quand les différences sont rares, de longues portions égales se parcourent sans remplir une surface quadratique. À l'inverse, sur une matrice dense, un débit en GCUPS\footnote{Milliards de mises à jour de cases par seconde. Pour une méthode dense, le débit nominal est \(mn/(10^9T)\). Les recalculs de halo ne figurent pas au numérateur.} permet de relier le temps de calcul aux dimensions. Les résultats des méthodes par bande, par bits ou par coût sont donc exprimés en secondes.

\section{Kokkos et calcul distribué}
Kokkos~\cite{kokkos} décrit un espace mémoire et un espace d'exécution. Le code utilise un tableau \texttt{View} pour la dernière occurrence et un second pour la DP. Un \texttt{parallel\_for} traite une ligne. Une même instance de l'espace d'exécution ordonne les lancements, puis une synchronisation finale précède le rapatriement du score. Le même code cible OpenMP sur EPYC et CUDA sur Hopper. La version native OpenMP garde une région parallèle autour des lignes. La version native CUDA peut aussi fusionner huit lignes par lancement. Un port de la récurrence n'assure donc pas les mêmes coûts de synchronisation.
Sur la paire de 20\,000 bases à huit cœurs, la médiane est de 27,90\,ms en OpenMP natif et de 110,22\,ms en Kokkos/OpenMP. Sur GH200, le même protocole donne 17,06\,ms en CUDA natif et 44,27\,ms en Kokkos/CUDA, avec 480\,024 octets de tampons explicites dans les deux cas. Ces mesures excluent l'initialisation du cadre d'exécution et suivent un appel d'échauffement. Le port Kokkos~4.5.01 a tourné sur GPU avec CUDA~12.6. L'essai avec CUDA~13 a échoué à la compilation de la bibliothèque.


\subsection{Deux nœuds complets}
Pour une paire, le rang \(r\) possède une bande verticale de colonnes. La dernière occurrence locale appartient à cette bande. Si elle manque, la position précédente du symbole dans les bandes de gauche est connue dès l'initialisation. Pour chaque ligne, chaque rang diffuse deux entiers, son bord droit et la valeur précédant sa dernière occurrence locale du symbole courant. Si cette occurrence est la première colonne du rang, son prédécesseur provient du bord droit du voisin. Un \texttt{Allgather} rassemble ces valeurs sur tous les rangs.
\begin{center}\begin{minipage}{.94\linewidth}
\textbf{Algorithme 5.} \textit{NvPD distribué par colonnes.} Les rangs exécutent les mêmes itérations.
\begin{algorithmic}[1]
\State Partager \(B\) en intervalles contigus, construire l'index local et repérer les occurrences antérieures
\For{chaque symbole \(c=A_i\)}
\State Rassembler les bords et prédécesseurs de tous les rangs pour la ligne \(i-1\)
\State Calculer en parallèle chaque colonne locale avec la récurrence de l'algorithme 1
\State Échanger les deux lignes locales
\EndFor
\State Communiquer la dernière case de la dernière bande
\end{algorithmic}
\end{minipage}\end{center}
Pour une matrice, le rang distribue des paires indépendantes à ses fils. Chaque fil traite sa paire sans collectif par ligne. Une réduction finale réunit le triangle supérieur. La figure ci-dessous distingue les points de communication.
\begin{figure}[H]\centering\begin{tikzpicture}[font=\footnotesize\rmfamily, align=center,
 boite/.style={fill=powderBlue!28,draw=powderBlue!55!black,rounded corners=1pt,minimum height=5mm,minimum width=21mm},
 lien/.style={-{Latex[length=1.4mm]},line width=.45pt}]
\foreach \x/\t in {0/{a) Une paire},7.1/{b) Lot de paires}}{
  \node[anchor=west,font=\footnotesize\bfseries\rmfamily] at (\x,1.21) {\t};}
\node[boite] (a) at (1.15,.45) {bande 0};
\node[boite] (b) at (4,.45) {bande 1};
\draw[lien] (a.east) -- (b.west);
\draw[lien] ([yshift=-1.2mm]b.west) -- ([yshift=-1.2mm]a.east);
\node[font=\scriptsize\rmfamily] at (2.58,-.18) {échange à chaque ligne};
\node[boite,fill=powderTeal!32] (c) at (8.25,.45) {lot 0};
\node[boite,fill=powderTeal!32] (d) at (11.10,.45) {lot 1};
\node[boite,fill=powderGold!48] (r) at (9.675,-.54) {réduction};
\draw[lien] (c.south) |- (r.west);\draw[lien] (d.south) |- (r.east);
\end{tikzpicture}

\caption{Deux stratégies MPI. Dans la paire distribuée (a), toutes les bandes échangent des bords à chaque ligne. Dans le lot indépendant (b), chaque distance reste locale jusqu'à la réduction finale.}
\end{figure}
\begin{table}[H]\centering
\caption{Même réservation exclusive, 24 rangs de huit fils par nœud. Les 192 puis 384 cœurs physiques distincts ont été observés. Médiane de trois maxima des temps de rang après échauffement, hors lecture du FASTA. L'efficacité à deux nœuds vaut \(S/2\).}
\small\begin{tabular}{lrrrr}\toprule
Charge & 1 nœud (s) & 2 nœuds (s) & \(S\) & Efficacité (\%)\\\midrule
\MpiRows
\bottomrule
\end{tabular}\end{table}
Les collectifs répétés rendent la paire plus lente sur deux nœuds aux deux formes essayées. Les 2\,016 paires de chaque matrice fournissent assez de calcul indépendant pour réduire le temps total.

\chapitre{Fenêtres homologues du chromosome humain}
Les fenêtres homologues comparent des segments correspondants de GRCh38.p14 et T2T-CHM13v2.0~\cite{t2t,chmreleases}. Sur le chromosome 21, les intervalles commencent en 34\,000\,000 et 32\,381\,795, en coordonnées de base zéro. Les ancres de 41 bases qui encadrent chaque intervalle sont uniques sur leur chromosome. La chaîne de correspondance UCSC place leurs extrémités dans le même ordre sur le brin positif. Les insertions et les délétions internes restent dans les fenêtres.
\chart{11_chromosome21.png}{Les fenêtres homologues et la copie du chromosome portant des insertions connues comparent les durées et les empreintes des méthodes.}{fig:humain}{bench/human.sbatch}
Sur la figure~\ref{fig:humain}, deux processus par méthode donnent les temps médians sur Grace et GH200. Pour les fenêtres de deux puis quatre millions de bases de référence, \(d=5\,599\) puis \(12\,482\). La portée \(R\) de l'index atteint 290 puis 622. Les deux calculs NvPD emploient donc un lancement par ligne. Leur temps passe de 18,350 à 90,891\,s, allocations, transferts et initialisation CUDA compris. Sur un cœur Grace, WFA2 demande 0,054 puis 0,348\,s et edlib 0,790 puis 2,991\,s.

Le contrôle entier ajoute 64 bases à une copie des 45\,090\,682 bases de CHM13. Le chemin construit coûte 64 et la différence de longueur impose \(d\ge64\). WFA2 obtient ce score en 11,34\,ms de calcul. Le processus prend 0,476\,s avec la lecture des deux chaînes.
\synthese{Sur les fenêtres homologues, un quadruplement approximatif du nombre de cases multiplie le temps du GPU par 4,95.}{À l'échelle d'un chromosome, la divergence entre versions d'assemblage peut varier de plusieurs ordres de grandeur.}

\chapitre{Comparaison des assemblages chromosomiques entiers}
Les versions v1.0 et v2.0 de CHM13 permettent de mesurer le coût d'une révision d'assemblage sur le chromosome 20. Les entrées \texttt{CP068258.1} et \texttt{CP068258.2} comptent 66\,210\,247 et 66\,210\,255 bases.\footnote{Voir l'\href{https://github.com/marbl/CHM13/blob/master/Earlier_assembly_releases_and_associated_data.md}{historique des assemblages CHM13 du consortium T2T}.} Le chromosome 20 porte déjà l'accession \texttt{CP068258.2} dans la version intermédiaire 1.1. Sur le chromosome 21, GRCh38 et CHM13 représentent deux références distinctes et comportent des bases ambiguës traitées comme des caractères ordinaires.

Une borne exploite l'histogramme des lettres. Pour \(u_c=\#_A(c)-\#_B(c)\) et \(x_+=\max(0,x)\), chaque édition efface au plus une unité d'excédent et une de déficit. Ainsi,
\[
 d(A,B)\ge\max\!\left(\sum_c(u_c)_+,\ \sum_c(-u_c)_+\right).
\]
Une borne supérieure suit un alignement global réalisable. Sur chaque bloc de la chaîne UCSC 23, il paie les désaccords. Dans chaque intervalle entre blocs, il choisit le moins coûteux des deux alignements ancrés à gauche ou à droite, puis paie l'écart de longueur. Les extrémités suivent la même règle. La somme des coûts des blocs et des intervalles donne le coût du chemin entier.
\chart{18_chromosomes_entiers.png}{Les durées sur le chromosome 20 et les bornes de distance sur le chromosome 21 comparent deux niveaux de divergence entre assemblages.}{fig:chrentiers}{bench/whole_chromosomes.py}
Les trois appels par méthode sur un cœur EPYC donnent \(d=30\) sur le chromosome 20. Le temps médian de calcul vaut 1,280\,s pour edlib et 0,092\,s pour WFA2. Avec lecture du FASTA, les processus prennent 2,679 et 1,484\,s. Sur le chromosome 21 de la figure~\ref{fig:chrentiers}, chaque processus atteint la limite murale de 120\,s. L'écart de longueur seul vaut 1\,619\,301 bases.
\synthese{Les fréquences des bases et le chemin construit encadrent la distance du chromosome 21 entre 6\,621\,364 et 16\,015\,052 éditions.}{Sur cette paire, la composition des séquences relève la borne de longueur d'un facteur 4,09.}


\chapitre{Géométrie des distances entre génomes de phage}
La reconstruction mesure dans quelle proportion les chemins d'un arbre reproduisent les distances globales de 64 séquences de phage \(\phi\)X174. Une origine différente sur le génome circulaire peut modifier la distance entre les chaînes déposées. Une même édition observée peut aussi résulter de plusieurs événements évolutifs.

La méthode UPGMA~\cite{upgma} joint les groupes \(U,V\) de distance moyenne minimale et place leur parent à la hauteur \(d(U,V)/2\). Pour un troisième groupe \(W\), la nouvelle distance moyenne est pondérée par les effectifs. Neighbor-Joining joint la paire qui minimise \(Q_{ij}\), où \(r\) est le nombre de groupes actifs et \(s_i=\sum_k d_{ik}\)~\cite{nj}.
\[
 d(U\cup V,W)=\frac{|U|d(U,W)+|V|d(V,W)}{|U|+|V|},\qquad
 Q_{ij}=(r-2)d_{ij}-s_i-s_j.
\]
Dans NJ, la première branche mesure \(d_{ij}/2+(s_i-s_j)/(2(r-2))\), la seconde complète \(d_{ij}\). À chaque fusion, les deux implémentations examinent toutes les paires actives. Pour \(N\) feuilles, elles utilisent \(O(N^3)\) opérations et \(O(N^2)\) distances stockées. UPGMA impose des feuilles équidistantes de la racine. Sur trois feuilles, les deux plus grandes distances doivent alors coïncider. L'additivité d'un arbre impose sur quatre feuilles l'égalité des deux plus grandes sommes parmi \(d_{ij}+d_{kl},d_{ik}+d_{jl},d_{il}+d_{jk}\).
\chart{16_phylogenie.png}{Les écarts aux conditions d'ultramétrie et d'additivité sont comparés aux résidus entre distances d'édition et longueurs des chemins.}{fig:phage}{scripts/phylo_fit.py}
La figure~\ref{fig:phage} porte sur les \(C(64,3)=41\,664\) triplets et \(C(64,4)=635\,376\) quartets. Les conditions sont satisfaites pour 1\,661 triplets et 245\,847 quartets. La racine de l'écart quadratique moyen vaut \(\sqrt{P^{-1}\sum_{i<j}(d_{ij}-t_{ij})^2}\), où \(t_{ij}\) est la longueur du chemin entre feuilles.
\synthese{Sur les 2\,016 paires, la racine de l'écart quadratique moyen vaut 14,89 éditions pour UPGMA et 3,56 pour NJ.}{La représentation des branches permet de localiser ces écarts dans les deux topologies.}

\chapitre{Topologie des arbres de phage}
Les phylogrammes de la figure~\ref{fig:arbres} utilisent la longueur cumulée des branches comme abscisse. Les échelles horizontales sont propres à chaque panneau. La longueur patristique se calcule par \(t_{ij}=\sum_{e\in\operatorname{chemin}(i,j)}\ell_e\). Le choix du point de dessin initial de NJ laisse ces sommes invariantes.
\begin{figure}[H]\centering
\resizebox{!}{.68\textheight}{\begin{tikzpicture}[x=1mm,y=-1mm,line width=.3pt]
\node[anchor=west,font=\small\bfseries,inner sep=0pt] at (0.000,-6) {a) UPGMA};
\draw[draw=powderBlue!50!black] (5.447,0.000) -- (34.880,0.000);
\draw[draw=powderBlue!50!black] (6.097,3.240) -- (34.880,3.240);
\draw[draw=powderBlue!50!black] (6.543,6.480) -- (34.880,6.480);
\draw[draw=powderBlue!50!black] (15.886,9.720) -- (34.880,9.720);
\draw[draw=powderBlue!50!black] (15.886,12.960) -- (34.880,12.960);
\draw[draw=powderBlue!50!black] (15.886,9.720) -- (15.886,12.960);
\draw[draw=powderBlue!50!black] (7.494,11.340) -- (15.886,11.340);
\draw[draw=powderBlue!50!black] (20.257,16.200) -- (34.880,16.200);
\draw[draw=powderBlue!50!black] (24.704,19.440) -- (34.880,19.440);
\draw[draw=powderBlue!50!black] (26.287,22.680) -- (34.880,22.680);
\draw[draw=powderBlue!50!black] (26.287,25.920) -- (34.880,25.920);
\draw[draw=powderBlue!50!black] (26.287,22.680) -- (26.287,25.920);
\draw[draw=powderBlue!50!black] (24.704,24.300) -- (26.287,24.300);
\draw[draw=powderBlue!50!black] (24.704,19.440) -- (24.704,24.300);
\draw[draw=powderBlue!50!black] (20.257,21.870) -- (24.704,21.870);
\draw[draw=powderBlue!50!black] (20.257,16.200) -- (20.257,21.870);
\draw[draw=powderBlue!50!black] (8.570,19.035) -- (20.257,19.035);
\draw[draw=powderBlue!50!black] (8.752,29.160) -- (34.880,29.160);
\draw[draw=powderBlue!50!black] (9.265,32.400) -- (34.880,32.400);
\draw[draw=powderBlue!50!black] (11.851,35.640) -- (34.880,35.640);
\draw[draw=powderBlue!50!black] (12.207,38.880) -- (34.880,38.880);
\draw[draw=powderBlue!50!black] (12.984,42.120) -- (34.880,42.120);
\draw[draw=powderBlue!50!black] (13.121,45.360) -- (34.880,45.360);
\draw[draw=powderBlue!50!black] (14.356,48.600) -- (34.880,48.600);
\draw[draw=powderBlue!50!black] (14.502,51.840) -- (34.880,51.840);
\draw[draw=powderBlue!50!black] (15.674,55.080) -- (34.880,55.080);
\draw[draw=powderBlue!50!black] (15.900,58.320) -- (34.880,58.320);
\draw[draw=powderBlue!50!black] (15.961,61.560) -- (34.880,61.560);
\draw[draw=powderBlue!50!black] (16.151,64.800) -- (34.880,64.800);
\draw[draw=powderBlue!50!black] (16.903,68.040) -- (34.880,68.040);
\draw[draw=powderBlue!50!black] (18.247,71.280) -- (34.880,71.280);
\draw[draw=powderBlue!50!black] (20.408,74.520) -- (34.880,74.520);
\draw[draw=powderBlue!50!black] (27.945,77.760) -- (34.880,77.760);
\draw[draw=powderBlue!50!black] (32.618,81.000) -- (34.880,81.000);
\draw[draw=powderBlue!50!black] (33.071,84.240) -- (34.880,84.240);
\draw[draw=powderBlue!50!black] (33.071,87.480) -- (34.880,87.480);
\draw[draw=powderBlue!50!black] (33.071,84.240) -- (33.071,87.480);
\draw[draw=powderBlue!50!black] (32.618,85.860) -- (33.071,85.860);
\draw[draw=powderBlue!50!black] (32.618,81.000) -- (32.618,85.860);
\draw[draw=powderBlue!50!black] (27.945,83.430) -- (32.618,83.430);
\draw[draw=powderBlue!50!black] (27.945,77.760) -- (27.945,83.430);
\draw[draw=powderBlue!50!black] (22.001,80.595) -- (27.945,80.595);
\draw[draw=powderBlue!50!black] (27.644,90.720) -- (34.880,90.720);
\draw[draw=powderBlue!50!black] (33.071,93.960) -- (34.880,93.960);
\draw[draw=powderBlue!50!black] (33.071,97.200) -- (34.880,97.200);
\draw[draw=powderBlue!50!black] (33.071,93.960) -- (33.071,97.200);
\draw[draw=powderBlue!50!black] (27.644,95.580) -- (33.071,95.580);
\draw[draw=powderBlue!50!black] (27.644,90.720) -- (27.644,95.580);
\draw[draw=powderBlue!50!black] (23.613,93.150) -- (27.644,93.150);
\draw[draw=powderBlue!50!black] (30.583,100.440) -- (34.880,100.440);
\draw[draw=powderBlue!50!black] (33.523,103.680) -- (34.880,103.680);
\draw[draw=powderBlue!50!black] (33.523,106.920) -- (34.880,106.920);
\draw[draw=powderBlue!50!black] (33.523,103.680) -- (33.523,106.920);
\draw[draw=powderBlue!50!black] (30.583,105.300) -- (33.523,105.300);
\draw[draw=powderBlue!50!black] (30.583,100.440) -- (30.583,105.300);
\draw[draw=powderBlue!50!black] (26.758,102.870) -- (30.583,102.870);
\draw[draw=powderBlue!50!black] (29.212,110.160) -- (34.880,110.160);
\draw[draw=powderBlue!50!black] (32.618,113.400) -- (34.880,113.400);
\draw[draw=powderBlue!50!black] (32.618,116.640) -- (34.880,116.640);
\draw[draw=powderBlue!50!black] (32.618,113.400) -- (32.618,116.640);
\draw[draw=powderBlue!50!black] (30.862,115.020) -- (32.618,115.020);
\draw[draw=powderBlue!50!black] (31.488,119.880) -- (34.880,119.880);
\draw[draw=powderBlue!50!black] (33.674,123.120) -- (34.880,123.120);
\draw[draw=powderBlue!50!black] (33.975,126.360) -- (34.880,126.360);
\draw[draw=powderBlue!50!black] (33.975,129.600) -- (34.880,129.600);
\draw[draw=powderBlue!50!black] (33.975,132.840) -- (34.880,132.840);
\draw[draw=powderBlue!50!black] (33.975,129.600) -- (33.975,132.840);
\draw[draw=powderBlue!50!black] (33.975,131.220) -- (33.975,131.220);
\draw[draw=powderBlue!50!black] (33.975,126.360) -- (33.975,131.220);
\draw[draw=powderBlue!50!black] (33.674,128.790) -- (33.975,128.790);
\draw[draw=powderBlue!50!black] (33.674,123.120) -- (33.674,128.790);
\draw[draw=powderBlue!50!black] (32.505,125.955) -- (33.674,125.955);
\draw[draw=powderBlue!50!black] (33.975,136.080) -- (34.880,136.080);
\draw[draw=powderBlue!50!black] (33.975,139.320) -- (34.880,139.320);
\draw[draw=powderBlue!50!black] (33.975,136.080) -- (33.975,139.320);
\draw[draw=powderBlue!50!black] (33.221,137.700) -- (33.975,137.700);
\draw[draw=powderBlue!50!black] (33.613,142.560) -- (34.880,142.560);
\draw[draw=powderBlue!50!black] (34.880,145.800) -- (34.880,145.800);
\draw[draw=powderBlue!50!black] (34.880,149.040) -- (34.880,149.040);
\draw[draw=powderBlue!50!black] (34.880,145.800) -- (34.880,149.040);
\draw[draw=powderBlue!50!black] (34.126,147.420) -- (34.880,147.420);
\draw[draw=powderBlue!50!black] (34.201,152.280) -- (34.880,152.280);
\draw[draw=powderBlue!50!black] (34.427,155.520) -- (34.880,155.520);
\draw[draw=powderBlue!50!black] (34.427,158.760) -- (34.880,158.760);
\draw[draw=powderBlue!50!black] (34.427,155.520) -- (34.427,158.760);
\draw[draw=powderBlue!50!black] (34.201,157.140) -- (34.427,157.140);
\draw[draw=powderBlue!50!black] (34.201,152.280) -- (34.201,157.140);
\draw[draw=powderBlue!50!black] (34.126,154.710) -- (34.201,154.710);
\draw[draw=powderBlue!50!black] (34.126,147.420) -- (34.126,154.710);
\draw[draw=powderBlue!50!black] (33.613,151.065) -- (34.126,151.065);
\draw[draw=powderBlue!50!black] (33.613,142.560) -- (33.613,151.065);
\draw[draw=powderBlue!50!black] (33.221,146.812) -- (33.613,146.812);
\draw[draw=powderBlue!50!black] (33.221,137.700) -- (33.221,146.812);
\draw[draw=powderBlue!50!black] (32.505,142.256) -- (33.221,142.256);
\draw[draw=powderBlue!50!black] (32.505,125.955) -- (32.505,142.256);
\draw[draw=powderBlue!50!black] (31.488,134.106) -- (32.505,134.106);
\draw[draw=powderBlue!50!black] (31.488,119.880) -- (31.488,134.106);
\draw[draw=powderBlue!50!black] (30.862,126.993) -- (31.488,126.993);
\draw[draw=powderBlue!50!black] (30.862,115.020) -- (30.862,126.993);
\draw[draw=powderBlue!50!black] (29.212,121.006) -- (30.862,121.006);
\draw[draw=powderBlue!50!black] (29.212,110.160) -- (29.212,121.006);
\draw[draw=powderBlue!50!black] (26.758,115.583) -- (29.212,115.583);
\draw[draw=powderBlue!50!black] (26.758,102.870) -- (26.758,115.583);
\draw[draw=powderBlue!50!black] (23.613,109.227) -- (26.758,109.227);
\draw[draw=powderBlue!50!black] (23.613,93.150) -- (23.613,109.227);
\draw[draw=powderBlue!50!black] (22.001,101.188) -- (23.613,101.188);
\draw[draw=powderBlue!50!black] (22.001,80.595) -- (22.001,101.188);
\draw[draw=powderBlue!50!black] (20.408,90.892) -- (22.001,90.892);
\draw[draw=powderBlue!50!black] (20.408,74.520) -- (20.408,90.892);
\draw[draw=powderBlue!50!black] (18.247,82.706) -- (20.408,82.706);
\draw[draw=powderBlue!50!black] (18.247,71.280) -- (18.247,82.706);
\draw[draw=powderBlue!50!black] (16.903,76.993) -- (18.247,76.993);
\draw[draw=powderBlue!50!black] (16.903,68.040) -- (16.903,76.993);
\draw[draw=powderBlue!50!black] (16.151,72.516) -- (16.903,72.516);
\draw[draw=powderBlue!50!black] (16.151,64.800) -- (16.151,72.516);
\draw[draw=powderBlue!50!black] (15.961,68.658) -- (16.151,68.658);
\draw[draw=powderBlue!50!black] (15.961,61.560) -- (15.961,68.658);
\draw[draw=powderBlue!50!black] (15.900,65.109) -- (15.961,65.109);
\draw[draw=powderBlue!50!black] (15.900,58.320) -- (15.900,65.109);
\draw[draw=powderBlue!50!black] (15.674,61.715) -- (15.900,61.715);
\draw[draw=powderBlue!50!black] (15.674,55.080) -- (15.674,61.715);
\draw[draw=powderBlue!50!black] (14.502,58.397) -- (15.674,58.397);
\draw[draw=powderBlue!50!black] (14.502,51.840) -- (14.502,58.397);
\draw[draw=powderBlue!50!black] (14.356,55.119) -- (14.502,55.119);
\draw[draw=powderBlue!50!black] (14.356,48.600) -- (14.356,55.119);
\draw[draw=powderBlue!50!black] (13.121,51.859) -- (14.356,51.859);
\draw[draw=powderBlue!50!black] (13.121,45.360) -- (13.121,51.859);
\draw[draw=powderBlue!50!black] (12.984,48.610) -- (13.121,48.610);
\draw[draw=powderBlue!50!black] (12.984,42.120) -- (12.984,48.610);
\draw[draw=powderBlue!50!black] (12.207,45.365) -- (12.984,45.365);
\draw[draw=powderBlue!50!black] (12.207,38.880) -- (12.207,45.365);
\draw[draw=powderBlue!50!black] (11.851,42.122) -- (12.207,42.122);
\draw[draw=powderBlue!50!black] (11.851,35.640) -- (11.851,42.122);
\draw[draw=powderBlue!50!black] (9.265,38.881) -- (11.851,38.881);
\draw[draw=powderBlue!50!black] (9.265,32.400) -- (9.265,38.881);
\draw[draw=powderBlue!50!black] (8.752,35.641) -- (9.265,35.641);
\draw[draw=powderBlue!50!black] (8.752,29.160) -- (8.752,35.641);
\draw[draw=powderBlue!50!black] (8.570,32.400) -- (8.752,32.400);
\draw[draw=powderBlue!50!black] (8.570,19.035) -- (8.570,32.400);
\draw[draw=powderBlue!50!black] (7.494,25.718) -- (8.570,25.718);
\draw[draw=powderBlue!50!black] (7.494,11.340) -- (7.494,25.718);
\draw[draw=powderBlue!50!black] (6.543,18.529) -- (7.494,18.529);
\draw[draw=powderBlue!50!black] (6.543,6.480) -- (6.543,18.529);
\draw[draw=powderBlue!50!black] (6.097,12.504) -- (6.543,12.504);
\draw[draw=powderBlue!50!black] (6.097,3.240) -- (6.097,12.504);
\draw[draw=powderBlue!50!black] (5.447,7.872) -- (6.097,7.872);
\draw[draw=powderBlue!50!black] (5.447,0.000) -- (5.447,7.872);
\draw[draw=powderBlue!50!black] (0.000,3.936) -- (5.447,3.936);
\draw[draw=powderBlue!50!black] (16.338,162.000) -- (34.880,162.000);
\draw[draw=powderBlue!50!black] (31.714,165.240) -- (34.880,165.240);
\draw[draw=powderBlue!50!black] (31.714,168.480) -- (34.880,168.480);
\draw[draw=powderBlue!50!black] (31.714,165.240) -- (31.714,168.480);
\draw[draw=powderBlue!50!black] (16.338,166.860) -- (31.714,166.860);
\draw[draw=powderBlue!50!black] (16.338,162.000) -- (16.338,166.860);
\draw[draw=powderBlue!50!black] (0.976,164.430) -- (16.338,164.430);
\draw[draw=powderBlue!50!black] (24.704,171.720) -- (34.880,171.720);
\draw[draw=powderBlue!50!black] (31.714,174.960) -- (34.880,174.960);
\draw[draw=powderBlue!50!black] (31.714,178.200) -- (34.880,178.200);
\draw[draw=powderBlue!50!black] (31.714,174.960) -- (31.714,178.200);
\draw[draw=powderBlue!50!black] (24.704,176.580) -- (31.714,176.580);
\draw[draw=powderBlue!50!black] (24.704,171.720) -- (24.704,176.580);
\draw[draw=powderBlue!50!black] (11.797,174.150) -- (24.704,174.150);
\draw[draw=powderBlue!50!black] (31.262,181.440) -- (34.880,181.440);
\draw[draw=powderBlue!50!black] (32.166,184.680) -- (34.880,184.680);
\draw[draw=powderBlue!50!black] (32.166,187.920) -- (34.880,187.920);
\draw[draw=powderBlue!50!black] (32.166,184.680) -- (32.166,187.920);
\draw[draw=powderBlue!50!black] (31.262,186.300) -- (32.166,186.300);
\draw[draw=powderBlue!50!black] (31.262,181.440) -- (31.262,186.300);
\draw[draw=powderBlue!50!black] (16.278,183.870) -- (31.262,183.870);
\draw[draw=powderBlue!50!black] (22.217,191.160) -- (34.880,191.160);
\draw[draw=powderBlue!50!black] (28.850,194.400) -- (34.880,194.400);
\draw[draw=powderBlue!50!black] (29.679,197.640) -- (34.880,197.640);
\draw[draw=powderBlue!50!black] (31.714,200.880) -- (34.880,200.880);
\draw[draw=powderBlue!50!black] (31.714,204.120) -- (34.880,204.120);
\draw[draw=powderBlue!50!black] (31.714,200.880) -- (31.714,204.120);
\draw[draw=powderBlue!50!black] (29.679,202.500) -- (31.714,202.500);
\draw[draw=powderBlue!50!black] (29.679,197.640) -- (29.679,202.500);
\draw[draw=powderBlue!50!black] (28.850,200.070) -- (29.679,200.070);
\draw[draw=powderBlue!50!black] (28.850,194.400) -- (28.850,200.070);
\draw[draw=powderBlue!50!black] (22.217,197.235) -- (28.850,197.235);
\draw[draw=powderBlue!50!black] (22.217,191.160) -- (22.217,197.235);
\draw[draw=powderBlue!50!black] (16.278,194.198) -- (22.217,194.198);
\draw[draw=powderBlue!50!black] (16.278,183.870) -- (16.278,194.198);
\draw[draw=powderBlue!50!black] (11.797,189.034) -- (16.278,189.034);
\draw[draw=powderBlue!50!black] (11.797,174.150) -- (11.797,189.034);
\draw[draw=powderBlue!50!black] (0.976,181.592) -- (11.797,181.592);
\draw[draw=powderBlue!50!black] (0.976,164.430) -- (0.976,181.592);
\draw[draw=powderBlue!50!black] (0.000,173.011) -- (0.976,173.011);
\draw[draw=powderBlue!50!black] (0.000,3.936) -- (0.000,173.011);
\draw[densely dotted,black!36] (34.880,0.000) -- (35.880,0.000);
\node[anchor=west,inner sep=0pt,font=\rmfamily\fontsize{7}{7}\selectfont] at (36.880,0.000) {EF380022.1};
\draw[densely dotted,black!36] (34.880,3.240) -- (35.880,3.240);
\node[anchor=west,inner sep=0pt,font=\rmfamily\fontsize{7}{7}\selectfont] at (36.880,3.240) {EF380026.1};
\draw[densely dotted,black!36] (34.880,6.480) -- (35.880,6.480);
\node[anchor=west,inner sep=0pt,font=\rmfamily\fontsize{7}{7}\selectfont] at (36.880,6.480) {EF380023.1};
\draw[densely dotted,black!36] (34.880,9.720) -- (35.880,9.720);
\node[anchor=west,inner sep=0pt,font=\rmfamily\fontsize{7}{7}\selectfont] at (36.880,9.720) {EF380024.1};
\draw[densely dotted,black!36] (34.880,12.960) -- (35.880,12.960);
\node[anchor=west,inner sep=0pt,font=\rmfamily\fontsize{7}{7}\selectfont] at (36.880,12.960) {EF380025.1};
\draw[densely dotted,black!36] (34.880,16.200) -- (35.880,16.200);
\node[anchor=west,inner sep=0pt,font=\rmfamily\fontsize{7}{7}\selectfont] at (36.880,16.200) {EF380029.1};
\draw[densely dotted,black!36] (34.880,19.440) -- (35.880,19.440);
\node[anchor=west,inner sep=0pt,font=\rmfamily\fontsize{7}{7}\selectfont] at (36.880,19.440) {EF380021.1};
\draw[densely dotted,black!36] (34.880,22.680) -- (35.880,22.680);
\node[anchor=west,inner sep=0pt,font=\rmfamily\fontsize{7}{7}\selectfont] at (36.880,22.680) {EF380028.1};
\draw[densely dotted,black!36] (34.880,25.920) -- (35.880,25.920);
\node[anchor=west,inner sep=0pt,font=\rmfamily\fontsize{7}{7}\selectfont] at (36.880,25.920) {EF380030.1};
\draw[densely dotted,black!36] (34.880,29.160) -- (35.880,29.160);
\node[anchor=west,inner sep=0pt,font=\rmfamily\fontsize{7}{7}\selectfont] at (36.880,29.160) {EF380027.1};
\draw[densely dotted,black!36] (34.880,32.400) -- (35.880,32.400);
\node[anchor=west,inner sep=0pt,font=\rmfamily\fontsize{7}{7}\selectfont] at (36.880,32.400) {EF380012.1};
\draw[densely dotted,black!36] (34.880,35.640) -- (35.880,35.640);
\node[anchor=west,inner sep=0pt,font=\rmfamily\fontsize{7}{7}\selectfont] at (36.880,35.640) {EF380013.1};
\draw[densely dotted,black!36] (34.880,38.880) -- (35.880,38.880);
\node[anchor=west,inner sep=0pt,font=\rmfamily\fontsize{7}{7}\selectfont] at (36.880,38.880) {EF380011.1};
\draw[densely dotted,black!36] (34.880,42.120) -- (35.880,42.120);
\node[anchor=west,inner sep=0pt,font=\rmfamily\fontsize{7}{7}\selectfont] at (36.880,42.120) {EF380014.1};
\draw[densely dotted,black!36] (34.880,45.360) -- (35.880,45.360);
\node[anchor=west,inner sep=0pt,font=\rmfamily\fontsize{7}{7}\selectfont] at (36.880,45.360) {EF380018.1};
\draw[densely dotted,black!36] (34.880,48.600) -- (35.880,48.600);
\node[anchor=west,inner sep=0pt,font=\rmfamily\fontsize{7}{7}\selectfont] at (36.880,48.600) {GQ153914.1};
\draw[densely dotted,black!36] (34.880,51.840) -- (35.880,51.840);
\node[anchor=west,inner sep=0pt,font=\rmfamily\fontsize{7}{7}\selectfont] at (36.880,51.840) {GQ153913.1};
\draw[densely dotted,black!36] (34.880,55.080) -- (35.880,55.080);
\node[anchor=west,inner sep=0pt,font=\rmfamily\fontsize{7}{7}\selectfont] at (36.880,55.080) {EF380010.1};
\draw[densely dotted,black!36] (34.880,58.320) -- (35.880,58.320);
\node[anchor=west,inner sep=0pt,font=\rmfamily\fontsize{7}{7}\selectfont] at (36.880,58.320) {EF380016.1};
\draw[densely dotted,black!36] (34.880,61.560) -- (35.880,61.560);
\node[anchor=west,inner sep=0pt,font=\rmfamily\fontsize{7}{7}\selectfont] at (36.880,61.560) {EF380019.1};
\draw[densely dotted,black!36] (34.880,64.800) -- (35.880,64.800);
\node[anchor=west,inner sep=0pt,font=\rmfamily\fontsize{7}{7}\selectfont] at (36.880,64.800) {EF380020.1};
\draw[densely dotted,black!36] (34.880,68.040) -- (35.880,68.040);
\node[anchor=west,inner sep=0pt,font=\rmfamily\fontsize{7}{7}\selectfont] at (36.880,68.040) {EF380015.1};
\draw[densely dotted,black!36] (34.880,71.280) -- (35.880,71.280);
\node[anchor=west,inner sep=0pt,font=\rmfamily\fontsize{7}{7}\selectfont] at (36.880,71.280) {EF380017.1};
\draw[densely dotted,black!36] (34.880,74.520) -- (35.880,74.520);
\node[anchor=west,inner sep=0pt,font=\rmfamily\fontsize{7}{7}\selectfont] at (36.880,74.520) {GQ153912.1};
\draw[densely dotted,black!36] (34.880,77.760) -- (35.880,77.760);
\node[anchor=west,inner sep=0pt,font=\rmfamily\fontsize{7}{7}\selectfont] at (36.880,77.760) {AF176027.1};
\draw[densely dotted,black!36] (34.880,81.000) -- (35.880,81.000);
\node[anchor=west,inner sep=0pt,font=\rmfamily\fontsize{7}{7}\selectfont] at (36.880,81.000) {AF176033.1};
\draw[densely dotted,black!36] (34.880,84.240) -- (35.880,84.240);
\node[anchor=west,inner sep=0pt,font=\rmfamily\fontsize{7}{7}\selectfont] at (36.880,84.240) {AF176029.1};
\draw[densely dotted,black!36] (34.880,87.480) -- (35.880,87.480);
\node[anchor=west,inner sep=0pt,font=\rmfamily\fontsize{7}{7}\selectfont] at (36.880,87.480) {AF176031.1};
\draw[densely dotted,black!36] (34.880,90.720) -- (35.880,90.720);
\node[anchor=west,inner sep=0pt,font=\rmfamily\fontsize{7}{7}\selectfont] at (36.880,90.720) {AF299313.1};
\draw[densely dotted,black!36] (34.880,93.960) -- (35.880,93.960);
\node[anchor=west,inner sep=0pt,font=\rmfamily\fontsize{7}{7}\selectfont] at (36.880,93.960) {AF299301.1};
\draw[densely dotted,black!36] (34.880,97.200) -- (35.880,97.200);
\node[anchor=west,inner sep=0pt,font=\rmfamily\fontsize{7}{7}\selectfont] at (36.880,97.200) {AF299305.1};
\draw[densely dotted,black!36] (34.880,100.440) -- (35.880,100.440);
\node[anchor=west,inner sep=0pt,font=\rmfamily\fontsize{7}{7}\selectfont] at (36.880,100.440) {EF380009.1};
\draw[densely dotted,black!36] (34.880,103.680) -- (35.880,103.680);
\node[anchor=west,inner sep=0pt,font=\rmfamily\fontsize{7}{7}\selectfont] at (36.880,103.680) {EF380031.1};
\draw[densely dotted,black!36] (34.880,106.920) -- (35.880,106.920);
\node[anchor=west,inner sep=0pt,font=\rmfamily\fontsize{7}{7}\selectfont] at (36.880,106.920) {EF380032.1};
\draw[densely dotted,black!36] (34.880,110.160) -- (35.880,110.160);
\node[anchor=west,inner sep=0pt,font=\rmfamily\fontsize{7}{7}\selectfont] at (36.880,110.160) {GQ153915.1};
\draw[densely dotted,black!36] (34.880,113.400) -- (35.880,113.400);
\node[anchor=west,inner sep=0pt,font=\rmfamily\fontsize{7}{7}\selectfont] at (36.880,113.400) {NC\_001422.1};
\draw[densely dotted,black!36] (34.880,116.640) -- (35.880,116.640);
\node[anchor=west,inner sep=0pt,font=\rmfamily\fontsize{7}{7}\selectfont] at (36.880,116.640) {CP004084.1};
\draw[densely dotted,black!36] (34.880,119.880) -- (35.880,119.880);
\node[anchor=west,inner sep=0pt,font=\rmfamily\fontsize{7}{7}\selectfont] at (36.880,119.880) {AF299302.1};
\draw[densely dotted,black!36] (34.880,123.120) -- (35.880,123.120);
\node[anchor=west,inner sep=0pt,font=\rmfamily\fontsize{7}{7}\selectfont] at (36.880,123.120) {HM753662.1};
\draw[densely dotted,black!36] (34.880,126.360) -- (35.880,126.360);
\node[anchor=west,inner sep=0pt,font=\rmfamily\fontsize{7}{7}\selectfont] at (36.880,126.360) {GU385905.1};
\draw[densely dotted,black!36] (34.880,129.600) -- (35.880,129.600);
\node[anchor=west,inner sep=0pt,font=\rmfamily\fontsize{7}{7}\selectfont] at (36.880,129.600) {AF176034.1};
\draw[densely dotted,black!36] (34.880,132.840) -- (35.880,132.840);
\node[anchor=west,inner sep=0pt,font=\rmfamily\fontsize{7}{7}\selectfont] at (36.880,132.840) {FJ849058.1};
\draw[densely dotted,black!36] (34.880,136.080) -- (35.880,136.080);
\node[anchor=west,inner sep=0pt,font=\rmfamily\fontsize{7}{7}\selectfont] at (36.880,136.080) {HM753664.1};
\draw[densely dotted,black!36] (34.880,139.320) -- (35.880,139.320);
\node[anchor=west,inner sep=0pt,font=\rmfamily\fontsize{7}{7}\selectfont] at (36.880,139.320) {HM753665.1};
\draw[densely dotted,black!36] (34.880,142.560) -- (35.880,142.560);
\node[anchor=west,inner sep=0pt,font=\rmfamily\fontsize{7}{7}\selectfont] at (36.880,142.560) {HM753658.1};
\draw[densely dotted,black!36] (34.880,145.800) -- (35.880,145.800);
\node[anchor=west,inner sep=0pt,font=\rmfamily\fontsize{7}{7}\selectfont] at (36.880,145.800) {HM753660.1};
\draw[densely dotted,black!36] (34.880,149.040) -- (35.880,149.040);
\node[anchor=west,inner sep=0pt,font=\rmfamily\fontsize{7}{7}\selectfont] at (36.880,149.040) {HM753661.1};
\draw[densely dotted,black!36] (34.880,152.280) -- (35.880,152.280);
\node[anchor=west,inner sep=0pt,font=\rmfamily\fontsize{7}{7}\selectfont] at (36.880,152.280) {HM753666.1};
\draw[densely dotted,black!36] (34.880,155.520) -- (35.880,155.520);
\node[anchor=west,inner sep=0pt,font=\rmfamily\fontsize{7}{7}\selectfont] at (36.880,155.520) {HM753659.1};
\draw[densely dotted,black!36] (34.880,158.760) -- (35.880,158.760);
\node[anchor=west,inner sep=0pt,font=\rmfamily\fontsize{7}{7}\selectfont] at (36.880,158.760) {HM753663.1};
\draw[densely dotted,black!36] (34.880,162.000) -- (35.880,162.000);
\node[anchor=west,inner sep=0pt,font=\rmfamily\fontsize{7}{7}\selectfont] at (36.880,162.000) {AF299307.1};
\draw[densely dotted,black!36] (34.880,165.240) -- (35.880,165.240);
\node[anchor=west,inner sep=0pt,font=\rmfamily\fontsize{7}{7}\selectfont] at (36.880,165.240) {AF299306.1};
\draw[densely dotted,black!36] (34.880,168.480) -- (35.880,168.480);
\node[anchor=west,inner sep=0pt,font=\rmfamily\fontsize{7}{7}\selectfont] at (36.880,168.480) {AF299308.1};
\draw[densely dotted,black!36] (34.880,171.720) -- (35.880,171.720);
\node[anchor=west,inner sep=0pt,font=\rmfamily\fontsize{7}{7}\selectfont] at (36.880,171.720) {AF299312.1};
\draw[densely dotted,black!36] (34.880,174.960) -- (35.880,174.960);
\node[anchor=west,inner sep=0pt,font=\rmfamily\fontsize{7}{7}\selectfont] at (36.880,174.960) {AF299303.1};
\draw[densely dotted,black!36] (34.880,178.200) -- (35.880,178.200);
\node[anchor=west,inner sep=0pt,font=\rmfamily\fontsize{7}{7}\selectfont] at (36.880,178.200) {AF299311.1};
\draw[densely dotted,black!36] (34.880,181.440) -- (35.880,181.440);
\node[anchor=west,inner sep=0pt,font=\rmfamily\fontsize{7}{7}\selectfont] at (36.880,181.440) {AF176032.1};
\draw[densely dotted,black!36] (34.880,184.680) -- (35.880,184.680);
\node[anchor=west,inner sep=0pt,font=\rmfamily\fontsize{7}{7}\selectfont] at (36.880,184.680) {AF176028.1};
\draw[densely dotted,black!36] (34.880,187.920) -- (35.880,187.920);
\node[anchor=west,inner sep=0pt,font=\rmfamily\fontsize{7}{7}\selectfont] at (36.880,187.920) {AF176030.1};
\draw[densely dotted,black!36] (34.880,191.160) -- (35.880,191.160);
\node[anchor=west,inner sep=0pt,font=\rmfamily\fontsize{7}{7}\selectfont] at (36.880,191.160) {AF299314.1};
\draw[densely dotted,black!36] (34.880,194.400) -- (35.880,194.400);
\node[anchor=west,inner sep=0pt,font=\rmfamily\fontsize{7}{7}\selectfont] at (36.880,194.400) {AF299309.1};
\draw[densely dotted,black!36] (34.880,197.640) -- (35.880,197.640);
\node[anchor=west,inner sep=0pt,font=\rmfamily\fontsize{7}{7}\selectfont] at (36.880,197.640) {AF299310.1};
\draw[densely dotted,black!36] (34.880,200.880) -- (35.880,200.880);
\node[anchor=west,inner sep=0pt,font=\rmfamily\fontsize{7}{7}\selectfont] at (36.880,200.880) {AF299300.1};
\draw[densely dotted,black!36] (34.880,204.120) -- (35.880,204.120);
\node[anchor=west,inner sep=0pt,font=\rmfamily\fontsize{7}{7}\selectfont] at (36.880,204.120) {AF299304.1};
\draw[line width=.6pt] (0.000,210) -- ++(9.045,0) node[midway,below=2pt,font=\scriptsize] {10 éditions};
\node[anchor=west,font=\small\bfseries,inner sep=0pt] at (62.880,-6) {b) Neighbor-Joining};
\draw[draw=powderTeal!50!black] (62.887,0.000) -- (62.887,0.000);
\draw[draw=powderTeal!50!black] (62.887,3.240) -- (63.791,3.240);
\draw[draw=powderTeal!50!black] (62.887,0.000) -- (62.887,3.240);
\draw[draw=powderTeal!50!black] (62.880,1.620) -- (62.887,1.620);
\draw[draw=powderTeal!50!black] (63.777,6.480) -- (63.792,6.480);
\draw[draw=powderTeal!50!black] (64.666,9.720) -- (65.571,9.720);
\draw[draw=powderTeal!50!black] (64.666,12.960) -- (65.571,12.960);
\draw[draw=powderTeal!50!black] (64.666,9.720) -- (64.666,12.960);
\draw[draw=powderTeal!50!black] (63.777,11.340) -- (64.666,11.340);
\draw[draw=powderTeal!50!black] (63.777,6.480) -- (63.777,11.340);
\draw[draw=powderTeal!50!black] (62.887,8.910) -- (63.777,8.910);
\draw[draw=powderTeal!50!black] (62.904,16.200) -- (64.691,16.200);
\draw[draw=powderTeal!50!black] (64.603,19.440) -- (65.372,19.440);
\draw[draw=powderTeal!50!black] (75.779,22.680) -- (78.780,22.680);
\draw[draw=powderTeal!50!black] (84.536,25.920) -- (85.554,25.920);
\draw[draw=powderTeal!50!black] (85.327,29.160) -- (86.952,29.160);
\draw[draw=powderTeal!50!black] (85.327,32.400) -- (88.225,32.400);
\draw[draw=powderTeal!50!black] (85.327,29.160) -- (85.327,32.400);
\draw[draw=powderTeal!50!black] (84.536,30.780) -- (85.327,30.780);
\draw[draw=powderTeal!50!black] (84.536,25.920) -- (84.536,30.780);
\draw[draw=powderTeal!50!black] (75.779,28.350) -- (84.536,28.350);
\draw[draw=powderTeal!50!black] (75.779,22.680) -- (75.779,28.350);
\draw[draw=powderTeal!50!black] (69.377,25.515) -- (75.779,25.515);
\draw[draw=powderTeal!50!black] (113.436,35.640) -- (114.751,35.640);
\draw[draw=powderTeal!50!black] (114.835,38.880) -- (117.586,38.880);
\draw[draw=powderTeal!50!black] (114.835,42.120) -- (119.320,42.120);
\draw[draw=powderTeal!50!black] (114.835,38.880) -- (114.835,42.120);
\draw[draw=powderTeal!50!black] (113.436,40.500) -- (114.835,40.500);
\draw[draw=powderTeal!50!black] (113.436,35.640) -- (113.436,40.500);
\draw[draw=powderTeal!50!black] (92.589,38.070) -- (113.436,38.070);
\draw[draw=powderTeal!50!black] (95.894,45.360) -- (109.861,45.360);
\draw[draw=powderTeal!50!black] (106.473,48.600) -- (130.888,48.600);
\draw[draw=powderTeal!50!black] (119.142,51.840) -- (120.402,51.840);
\draw[draw=powderTeal!50!black] (120.595,55.080) -- (125.992,55.080);
\draw[draw=powderTeal!50!black] (120.595,58.320) -- (136.000,58.320);
\draw[draw=powderTeal!50!black] (120.595,55.080) -- (120.595,58.320);
\draw[draw=powderTeal!50!black] (119.142,56.700) -- (120.595,56.700);
\draw[draw=powderTeal!50!black] (119.142,51.840) -- (119.142,56.700);
\draw[draw=powderTeal!50!black] (106.473,54.270) -- (119.142,54.270);
\draw[draw=powderTeal!50!black] (106.473,48.600) -- (106.473,54.270);
\draw[draw=powderTeal!50!black] (96.681,51.435) -- (106.473,51.435);
\draw[draw=powderTeal!50!black] (99.213,61.560) -- (100.225,61.560);
\draw[draw=powderTeal!50!black] (104.532,64.800) -- (104.768,64.800);
\draw[draw=powderTeal!50!black] (105.200,68.040) -- (112.521,68.040);
\draw[draw=powderTeal!50!black] (105.200,71.280) -- (111.447,71.280);
\draw[draw=powderTeal!50!black] (105.200,68.040) -- (105.200,71.280);
\draw[draw=powderTeal!50!black] (104.532,69.660) -- (105.200,69.660);
\draw[draw=powderTeal!50!black] (104.532,64.800) -- (104.532,69.660);
\draw[draw=powderTeal!50!black] (99.213,67.230) -- (104.532,67.230);
\draw[draw=powderTeal!50!black] (99.213,61.560) -- (99.213,67.230);
\draw[draw=powderTeal!50!black] (96.681,64.395) -- (99.213,64.395);
\draw[draw=powderTeal!50!black] (96.681,51.435) -- (96.681,64.395);
\draw[draw=powderTeal!50!black] (95.894,57.915) -- (96.681,57.915);
\draw[draw=powderTeal!50!black] (95.894,45.360) -- (95.894,57.915);
\draw[draw=powderTeal!50!black] (92.589,51.638) -- (95.894,51.638);
\draw[draw=powderTeal!50!black] (92.589,38.070) -- (92.589,51.638);
\draw[draw=powderTeal!50!black] (71.782,44.854) -- (92.589,44.854);
\draw[draw=powderTeal!50!black] (81.948,74.520) -- (91.095,74.520);
\draw[draw=powderTeal!50!black] (104.004,77.760) -- (104.106,77.760);
\draw[draw=powderTeal!50!black] (104.004,81.000) -- (110.234,81.000);
\draw[draw=powderTeal!50!black] (104.004,77.760) -- (104.004,81.000);
\draw[draw=powderTeal!50!black] (81.948,79.380) -- (104.004,79.380);
\draw[draw=powderTeal!50!black] (81.948,74.520) -- (81.948,79.380);
\draw[draw=powderTeal!50!black] (77.017,76.950) -- (81.948,76.950);
\draw[draw=powderTeal!50!black] (78.079,84.240) -- (78.155,84.240);
\draw[draw=powderTeal!50!black] (78.079,87.480) -- (81.620,87.480);
\draw[draw=powderTeal!50!black] (78.079,84.240) -- (78.079,87.480);
\draw[draw=powderTeal!50!black] (77.017,85.860) -- (78.079,85.860);
\draw[draw=powderTeal!50!black] (77.017,76.950) -- (77.017,85.860);
\draw[draw=powderTeal!50!black] (71.782,81.405) -- (77.017,81.405);
\draw[draw=powderTeal!50!black] (71.782,44.854) -- (71.782,81.405);
\draw[draw=powderTeal!50!black] (69.377,63.129) -- (71.782,63.129);
\draw[draw=powderTeal!50!black] (69.377,25.515) -- (69.377,63.129);
\draw[draw=powderTeal!50!black] (67.071,44.322) -- (69.377,44.322);
\draw[draw=powderTeal!50!black] (67.390,90.720) -- (69.264,90.720);
\draw[draw=powderTeal!50!black] (73.658,93.960) -- (73.752,93.960);
\draw[draw=powderTeal!50!black] (74.242,97.200) -- (95.196,97.200);
\draw[draw=powderTeal!50!black] (74.995,100.440) -- (87.119,100.440);
\draw[draw=powderTeal!50!black] (74.995,103.680) -- (95.432,103.680);
\draw[draw=powderTeal!50!black] (74.995,100.440) -- (74.995,103.680);
\draw[draw=powderTeal!50!black] (74.242,102.060) -- (74.995,102.060);
\draw[draw=powderTeal!50!black] (74.242,97.200) -- (74.242,102.060);
\draw[draw=powderTeal!50!black] (73.658,99.630) -- (74.242,99.630);
\draw[draw=powderTeal!50!black] (73.658,93.960) -- (73.658,99.630);
\draw[draw=powderTeal!50!black] (67.390,96.795) -- (73.658,96.795);
\draw[draw=powderTeal!50!black] (67.390,90.720) -- (67.390,96.795);
\draw[draw=powderTeal!50!black] (67.071,93.757) -- (67.390,93.757);
\draw[draw=powderTeal!50!black] (67.071,44.322) -- (67.071,93.757);
\draw[draw=powderTeal!50!black] (65.434,69.040) -- (67.071,69.040);
\draw[draw=powderTeal!50!black] (65.682,106.920) -- (65.845,106.920);
\draw[draw=powderTeal!50!black] (66.007,110.160) -- (67.078,110.160);
\draw[draw=powderTeal!50!black] (66.617,113.400) -- (67.541,113.400);
\draw[draw=powderTeal!50!black] (72.290,116.640) -- (73.039,116.640);
\draw[draw=powderTeal!50!black] (89.214,119.880) -- (105.073,119.880);
\draw[draw=powderTeal!50!black] (89.214,123.120) -- (111.343,123.120);
\draw[draw=powderTeal!50!black] (89.214,119.880) -- (89.214,123.120);
\draw[draw=powderTeal!50!black] (73.589,121.500) -- (89.214,121.500);
\draw[draw=powderTeal!50!black] (76.149,126.360) -- (96.536,126.360);
\draw[draw=powderTeal!50!black] (78.691,129.600) -- (96.047,129.600);
\draw[draw=powderTeal!50!black] (92.088,132.840) -- (94.326,132.840);
\draw[draw=powderTeal!50!black] (99.798,136.080) -- (118.515,136.080);
\draw[draw=powderTeal!50!black] (101.885,139.320) -- (110.257,139.320);
\draw[draw=powderTeal!50!black] (101.885,142.560) -- (110.698,142.560);
\draw[draw=powderTeal!50!black] (101.885,139.320) -- (101.885,142.560);
\draw[draw=powderTeal!50!black] (99.798,140.940) -- (101.885,140.940);
\draw[draw=powderTeal!50!black] (99.798,136.080) -- (99.798,140.940);
\draw[draw=powderTeal!50!black] (92.088,138.510) -- (99.798,138.510);
\draw[draw=powderTeal!50!black] (92.088,132.840) -- (92.088,138.510);
\draw[draw=powderTeal!50!black] (78.691,135.675) -- (92.088,135.675);
\draw[draw=powderTeal!50!black] (78.691,129.600) -- (78.691,135.675);
\draw[draw=powderTeal!50!black] (77.525,132.638) -- (78.691,132.638);
\draw[draw=powderTeal!50!black] (78.055,145.800) -- (100.571,145.800);
\draw[draw=powderTeal!50!black] (78.055,149.040) -- (100.762,149.040);
\draw[draw=powderTeal!50!black] (78.055,145.800) -- (78.055,149.040);
\draw[draw=powderTeal!50!black] (77.525,147.420) -- (78.055,147.420);
\draw[draw=powderTeal!50!black] (77.525,132.638) -- (77.525,147.420);
\draw[draw=powderTeal!50!black] (76.149,140.029) -- (77.525,140.029);
\draw[draw=powderTeal!50!black] (76.149,126.360) -- (76.149,140.029);
\draw[draw=powderTeal!50!black] (75.229,133.194) -- (76.149,133.194);
\draw[draw=powderTeal!50!black] (75.246,152.280) -- (94.723,152.280);
\draw[draw=powderTeal!50!black] (76.566,155.520) -- (96.476,155.520);
\draw[draw=powderTeal!50!black] (76.566,158.760) -- (96.453,158.760);
\draw[draw=powderTeal!50!black] (76.566,155.520) -- (76.566,158.760);
\draw[draw=powderTeal!50!black] (75.493,157.140) -- (76.566,157.140);
\draw[draw=powderTeal!50!black] (77.460,162.000) -- (102.542,162.000);
\draw[draw=powderTeal!50!black] (77.460,165.240) -- (92.175,165.240);
\draw[draw=powderTeal!50!black] (77.460,162.000) -- (77.460,165.240);
\draw[draw=powderTeal!50!black] (75.911,163.620) -- (77.460,163.620);
\draw[draw=powderTeal!50!black] (80.957,168.480) -- (113.848,168.480);
\draw[draw=powderTeal!50!black] (81.530,171.720) -- (112.307,171.720);
\draw[draw=powderTeal!50!black] (82.410,174.960) -- (111.363,174.960);
\draw[draw=powderTeal!50!black] (82.410,178.200) -- (108.630,178.200);
\draw[draw=powderTeal!50!black] (82.410,174.960) -- (82.410,178.200);
\draw[draw=powderTeal!50!black] (81.530,176.580) -- (82.410,176.580);
\draw[draw=powderTeal!50!black] (81.530,171.720) -- (81.530,176.580);
\draw[draw=powderTeal!50!black] (80.957,174.150) -- (81.530,174.150);
\draw[draw=powderTeal!50!black] (80.957,168.480) -- (80.957,174.150);
\draw[draw=powderTeal!50!black] (79.830,171.315) -- (80.957,171.315);
\draw[draw=powderTeal!50!black] (80.513,181.440) -- (82.312,181.440);
\draw[draw=powderTeal!50!black] (80.513,184.680) -- (81.427,184.680);
\draw[draw=powderTeal!50!black] (80.513,181.440) -- (80.513,184.680);
\draw[draw=powderTeal!50!black] (79.830,183.060) -- (80.513,183.060);
\draw[draw=powderTeal!50!black] (79.830,171.315) -- (79.830,183.060);
\draw[draw=powderTeal!50!black] (76.963,177.188) -- (79.830,177.188);
\draw[draw=powderTeal!50!black] (78.221,187.920) -- (102.157,187.920);
\draw[draw=powderTeal!50!black] (78.221,191.160) -- (107.649,191.160);
\draw[draw=powderTeal!50!black] (78.221,187.920) -- (78.221,191.160);
\draw[draw=powderTeal!50!black] (76.963,189.540) -- (78.221,189.540);
\draw[draw=powderTeal!50!black] (76.963,177.188) -- (76.963,189.540);
\draw[draw=powderTeal!50!black] (75.911,183.364) -- (76.963,183.364);
\draw[draw=powderTeal!50!black] (75.911,163.620) -- (75.911,183.364);
\draw[draw=powderTeal!50!black] (75.493,173.492) -- (75.911,173.492);
\draw[draw=powderTeal!50!black] (75.493,157.140) -- (75.493,173.492);
\draw[draw=powderTeal!50!black] (75.246,165.316) -- (75.493,165.316);
\draw[draw=powderTeal!50!black] (75.246,152.280) -- (75.246,165.316);
\draw[draw=powderTeal!50!black] (75.229,158.798) -- (75.246,158.798);
\draw[draw=powderTeal!50!black] (75.229,133.194) -- (75.229,158.798);
\draw[draw=powderTeal!50!black] (73.589,145.996) -- (75.229,145.996);
\draw[draw=powderTeal!50!black] (73.589,121.500) -- (73.589,145.996);
\draw[draw=powderTeal!50!black] (72.290,133.748) -- (73.589,133.748);
\draw[draw=powderTeal!50!black] (72.290,116.640) -- (72.290,133.748);
\draw[draw=powderTeal!50!black] (67.050,125.194) -- (72.290,125.194);
\draw[draw=powderTeal!50!black] (68.593,194.400) -- (70.459,194.400);
\draw[draw=powderTeal!50!black] (68.593,197.640) -- (71.250,197.640);
\draw[draw=powderTeal!50!black] (68.593,194.400) -- (68.593,197.640);
\draw[draw=powderTeal!50!black] (67.050,196.020) -- (68.593,196.020);
\draw[draw=powderTeal!50!black] (67.050,125.194) -- (67.050,196.020);
\draw[draw=powderTeal!50!black] (66.617,160.607) -- (67.050,160.607);
\draw[draw=powderTeal!50!black] (66.617,113.400) -- (66.617,160.607);
\draw[draw=powderTeal!50!black] (66.007,137.004) -- (66.617,137.004);
\draw[draw=powderTeal!50!black] (66.007,110.160) -- (66.007,137.004);
\draw[draw=powderTeal!50!black] (65.682,123.582) -- (66.007,123.582);
\draw[draw=powderTeal!50!black] (65.682,106.920) -- (65.682,123.582);
\draw[draw=powderTeal!50!black] (65.434,115.251) -- (65.682,115.251);
\draw[draw=powderTeal!50!black] (65.434,69.040) -- (65.434,115.251);
\draw[draw=powderTeal!50!black] (64.603,92.145) -- (65.434,92.145);
\draw[draw=powderTeal!50!black] (64.603,19.440) -- (64.603,92.145);
\draw[draw=powderTeal!50!black] (63.373,55.793) -- (64.603,55.793);
\draw[draw=powderTeal!50!black] (63.798,200.880) -- (63.798,200.880);
\draw[draw=powderTeal!50!black] (63.798,204.120) -- (63.798,204.120);
\draw[draw=powderTeal!50!black] (63.798,200.880) -- (63.798,204.120);
\draw[draw=powderTeal!50!black] (63.373,202.500) -- (63.798,202.500);
\draw[draw=powderTeal!50!black] (63.373,55.793) -- (63.373,202.500);
\draw[draw=powderTeal!50!black] (62.904,129.146) -- (63.373,129.146);
\draw[draw=powderTeal!50!black] (62.904,16.200) -- (62.904,129.146);
\draw[draw=powderTeal!50!black] (62.887,72.673) -- (62.904,72.673);
\draw[draw=powderTeal!50!black] (62.887,8.910) -- (62.887,72.673);
\draw[draw=powderTeal!50!black] (62.880,40.792) -- (62.887,40.792);
\draw[draw=powderTeal!50!black] (62.880,1.620) -- (62.880,40.792);
\draw[densely dotted,black!36] (62.887,0.000) -- (137.000,0.000);
\node[anchor=west,inner sep=0pt,font=\rmfamily\fontsize{7}{7}\selectfont] at (138.000,0.000) {HM753659.1};
\draw[densely dotted,black!36] (63.791,3.240) -- (137.000,3.240);
\node[anchor=west,inner sep=0pt,font=\rmfamily\fontsize{7}{7}\selectfont] at (138.000,3.240) {HM753663.1};
\draw[densely dotted,black!36] (63.792,6.480) -- (137.000,6.480);
\node[anchor=west,inner sep=0pt,font=\rmfamily\fontsize{7}{7}\selectfont] at (138.000,6.480) {HM753666.1};
\draw[densely dotted,black!36] (65.571,9.720) -- (137.000,9.720);
\node[anchor=west,inner sep=0pt,font=\rmfamily\fontsize{7}{7}\selectfont] at (138.000,9.720) {HM753664.1};
\draw[densely dotted,black!36] (65.571,12.960) -- (137.000,12.960);
\node[anchor=west,inner sep=0pt,font=\rmfamily\fontsize{7}{7}\selectfont] at (138.000,12.960) {HM753665.1};
\draw[densely dotted,black!36] (64.691,16.200) -- (137.000,16.200);
\node[anchor=west,inner sep=0pt,font=\rmfamily\fontsize{7}{7}\selectfont] at (138.000,16.200) {HM753658.1};
\draw[densely dotted,black!36] (65.372,19.440) -- (137.000,19.440);
\node[anchor=west,inner sep=0pt,font=\rmfamily\fontsize{7}{7}\selectfont] at (138.000,19.440) {HM753662.1};
\draw[densely dotted,black!36] (78.780,22.680) -- (137.000,22.680);
\node[anchor=west,inner sep=0pt,font=\rmfamily\fontsize{7}{7}\selectfont] at (138.000,22.680) {AF176027.1};
\draw[densely dotted,black!36] (85.554,25.920) -- (137.000,25.920);
\node[anchor=west,inner sep=0pt,font=\rmfamily\fontsize{7}{7}\selectfont] at (138.000,25.920) {AF176029.1};
\draw[densely dotted,black!36] (86.952,29.160) -- (137.000,29.160);
\node[anchor=west,inner sep=0pt,font=\rmfamily\fontsize{7}{7}\selectfont] at (138.000,29.160) {AF176031.1};
\draw[densely dotted,black!36] (88.225,32.400) -- (137.000,32.400);
\node[anchor=west,inner sep=0pt,font=\rmfamily\fontsize{7}{7}\selectfont] at (138.000,32.400) {AF176033.1};
\draw[densely dotted,black!36] (114.751,35.640) -- (137.000,35.640);
\node[anchor=west,inner sep=0pt,font=\rmfamily\fontsize{7}{7}\selectfont] at (138.000,35.640) {AF176028.1};
\draw[densely dotted,black!36] (117.586,38.880) -- (137.000,38.880);
\node[anchor=west,inner sep=0pt,font=\rmfamily\fontsize{7}{7}\selectfont] at (138.000,38.880) {AF176030.1};
\draw[densely dotted,black!36] (119.320,42.120) -- (137.000,42.120);
\node[anchor=west,inner sep=0pt,font=\rmfamily\fontsize{7}{7}\selectfont] at (138.000,42.120) {AF176032.1};
\draw[densely dotted,black!36] (109.861,45.360) -- (137.000,45.360);
\node[anchor=west,inner sep=0pt,font=\rmfamily\fontsize{7}{7}\selectfont] at (138.000,45.360) {AF299314.1};
\draw[densely dotted,black!36] (130.888,48.600) -- (137.000,48.600);
\node[anchor=west,inner sep=0pt,font=\rmfamily\fontsize{7}{7}\selectfont] at (138.000,48.600) {AF299307.1};
\draw[densely dotted,black!36] (120.402,51.840) -- (137.000,51.840);
\node[anchor=west,inner sep=0pt,font=\rmfamily\fontsize{7}{7}\selectfont] at (138.000,51.840) {AF299303.1};
\draw[densely dotted,black!36] (125.992,55.080) -- (137.000,55.080);
\node[anchor=west,inner sep=0pt,font=\rmfamily\fontsize{7}{7}\selectfont] at (138.000,55.080) {AF299311.1};
\draw[densely dotted,black!36] (136.000,58.320) -- (137.000,58.320);
\node[anchor=west,inner sep=0pt,font=\rmfamily\fontsize{7}{7}\selectfont] at (138.000,58.320) {AF299312.1};
\draw[densely dotted,black!36] (100.225,61.560) -- (137.000,61.560);
\node[anchor=west,inner sep=0pt,font=\rmfamily\fontsize{7}{7}\selectfont] at (138.000,61.560) {AF299300.1};
\draw[densely dotted,black!36] (104.768,64.800) -- (137.000,64.800);
\node[anchor=west,inner sep=0pt,font=\rmfamily\fontsize{7}{7}\selectfont] at (138.000,64.800) {AF299304.1};
\draw[densely dotted,black!36] (112.521,68.040) -- (137.000,68.040);
\node[anchor=west,inner sep=0pt,font=\rmfamily\fontsize{7}{7}\selectfont] at (138.000,68.040) {AF299309.1};
\draw[densely dotted,black!36] (111.447,71.280) -- (137.000,71.280);
\node[anchor=west,inner sep=0pt,font=\rmfamily\fontsize{7}{7}\selectfont] at (138.000,71.280) {AF299310.1};
\draw[densely dotted,black!36] (91.095,74.520) -- (137.000,74.520);
\node[anchor=west,inner sep=0pt,font=\rmfamily\fontsize{7}{7}\selectfont] at (138.000,74.520) {AF299313.1};
\draw[densely dotted,black!36] (104.106,77.760) -- (137.000,77.760);
\node[anchor=west,inner sep=0pt,font=\rmfamily\fontsize{7}{7}\selectfont] at (138.000,77.760) {AF299306.1};
\draw[densely dotted,black!36] (110.234,81.000) -- (137.000,81.000);
\node[anchor=west,inner sep=0pt,font=\rmfamily\fontsize{7}{7}\selectfont] at (138.000,81.000) {AF299308.1};
\draw[densely dotted,black!36] (78.155,84.240) -- (137.000,84.240);
\node[anchor=west,inner sep=0pt,font=\rmfamily\fontsize{7}{7}\selectfont] at (138.000,84.240) {AF299301.1};
\draw[densely dotted,black!36] (81.620,87.480) -- (137.000,87.480);
\node[anchor=west,inner sep=0pt,font=\rmfamily\fontsize{7}{7}\selectfont] at (138.000,87.480) {AF299305.1};
\draw[densely dotted,black!36] (69.264,90.720) -- (137.000,90.720);
\node[anchor=west,inner sep=0pt,font=\rmfamily\fontsize{7}{7}\selectfont] at (138.000,90.720) {AF299302.1};
\draw[densely dotted,black!36] (73.752,93.960) -- (137.000,93.960);
\node[anchor=west,inner sep=0pt,font=\rmfamily\fontsize{7}{7}\selectfont] at (138.000,93.960) {GQ153915.1};
\draw[densely dotted,black!36] (95.196,97.200) -- (137.000,97.200);
\node[anchor=west,inner sep=0pt,font=\rmfamily\fontsize{7}{7}\selectfont] at (138.000,97.200) {GQ153913.1};
\draw[densely dotted,black!36] (87.119,100.440) -- (137.000,100.440);
\node[anchor=west,inner sep=0pt,font=\rmfamily\fontsize{7}{7}\selectfont] at (138.000,100.440) {GQ153912.1};
\draw[densely dotted,black!36] (95.432,103.680) -- (137.000,103.680);
\node[anchor=west,inner sep=0pt,font=\rmfamily\fontsize{7}{7}\selectfont] at (138.000,103.680) {GQ153914.1};
\draw[densely dotted,black!36] (65.845,106.920) -- (137.000,106.920);
\node[anchor=west,inner sep=0pt,font=\rmfamily\fontsize{7}{7}\selectfont] at (138.000,106.920) {AF176034.1};
\draw[densely dotted,black!36] (67.078,110.160) -- (137.000,110.160);
\node[anchor=west,inner sep=0pt,font=\rmfamily\fontsize{7}{7}\selectfont] at (138.000,110.160) {FJ849058.1};
\draw[densely dotted,black!36] (67.541,113.400) -- (137.000,113.400);
\node[anchor=west,inner sep=0pt,font=\rmfamily\fontsize{7}{7}\selectfont] at (138.000,113.400) {GU385905.1};
\draw[densely dotted,black!36] (73.039,116.640) -- (137.000,116.640);
\node[anchor=west,inner sep=0pt,font=\rmfamily\fontsize{7}{7}\selectfont] at (138.000,116.640) {EF380009.1};
\draw[densely dotted,black!36] (105.073,119.880) -- (137.000,119.880);
\node[anchor=west,inner sep=0pt,font=\rmfamily\fontsize{7}{7}\selectfont] at (138.000,119.880) {EF380024.1};
\draw[densely dotted,black!36] (111.343,123.120) -- (137.000,123.120);
\node[anchor=west,inner sep=0pt,font=\rmfamily\fontsize{7}{7}\selectfont] at (138.000,123.120) {EF380025.1};
\draw[densely dotted,black!36] (96.536,126.360) -- (137.000,126.360);
\node[anchor=west,inner sep=0pt,font=\rmfamily\fontsize{7}{7}\selectfont] at (138.000,126.360) {EF380010.1};
\draw[densely dotted,black!36] (96.047,129.600) -- (137.000,129.600);
\node[anchor=west,inner sep=0pt,font=\rmfamily\fontsize{7}{7}\selectfont] at (138.000,129.600) {EF380020.1};
\draw[densely dotted,black!36] (94.326,132.840) -- (137.000,132.840);
\node[anchor=west,inner sep=0pt,font=\rmfamily\fontsize{7}{7}\selectfont] at (138.000,132.840) {EF380021.1};
\draw[densely dotted,black!36] (118.515,136.080) -- (137.000,136.080);
\node[anchor=west,inner sep=0pt,font=\rmfamily\fontsize{7}{7}\selectfont] at (138.000,136.080) {EF380029.1};
\draw[densely dotted,black!36] (110.257,139.320) -- (137.000,139.320);
\node[anchor=west,inner sep=0pt,font=\rmfamily\fontsize{7}{7}\selectfont] at (138.000,139.320) {EF380028.1};
\draw[densely dotted,black!36] (110.698,142.560) -- (137.000,142.560);
\node[anchor=west,inner sep=0pt,font=\rmfamily\fontsize{7}{7}\selectfont] at (138.000,142.560) {EF380030.1};
\draw[densely dotted,black!36] (100.571,145.800) -- (137.000,145.800);
\node[anchor=west,inner sep=0pt,font=\rmfamily\fontsize{7}{7}\selectfont] at (138.000,145.800) {EF380014.1};
\draw[densely dotted,black!36] (100.762,149.040) -- (137.000,149.040);
\node[anchor=west,inner sep=0pt,font=\rmfamily\fontsize{7}{7}\selectfont] at (138.000,149.040) {EF380018.1};
\draw[densely dotted,black!36] (94.723,152.280) -- (137.000,152.280);
\node[anchor=west,inner sep=0pt,font=\rmfamily\fontsize{7}{7}\selectfont] at (138.000,152.280) {EF380015.1};
\draw[densely dotted,black!36] (96.476,155.520) -- (137.000,155.520);
\node[anchor=west,inner sep=0pt,font=\rmfamily\fontsize{7}{7}\selectfont] at (138.000,155.520) {EF380016.1};
\draw[densely dotted,black!36] (96.453,158.760) -- (137.000,158.760);
\node[anchor=west,inner sep=0pt,font=\rmfamily\fontsize{7}{7}\selectfont] at (138.000,158.760) {EF380019.1};
\draw[densely dotted,black!36] (102.542,162.000) -- (137.000,162.000);
\node[anchor=west,inner sep=0pt,font=\rmfamily\fontsize{7}{7}\selectfont] at (138.000,162.000) {EF380013.1};
\draw[densely dotted,black!36] (92.175,165.240) -- (137.000,165.240);
\node[anchor=west,inner sep=0pt,font=\rmfamily\fontsize{7}{7}\selectfont] at (138.000,165.240) {EF380017.1};
\draw[densely dotted,black!36] (113.848,168.480) -- (137.000,168.480);
\node[anchor=west,inner sep=0pt,font=\rmfamily\fontsize{7}{7}\selectfont] at (138.000,168.480) {EF380022.1};
\draw[densely dotted,black!36] (112.307,171.720) -- (137.000,171.720);
\node[anchor=west,inner sep=0pt,font=\rmfamily\fontsize{7}{7}\selectfont] at (138.000,171.720) {EF380026.1};
\draw[densely dotted,black!36] (111.363,174.960) -- (137.000,174.960);
\node[anchor=west,inner sep=0pt,font=\rmfamily\fontsize{7}{7}\selectfont] at (138.000,174.960) {EF380023.1};
\draw[densely dotted,black!36] (108.630,178.200) -- (137.000,178.200);
\node[anchor=west,inner sep=0pt,font=\rmfamily\fontsize{7}{7}\selectfont] at (138.000,178.200) {EF380027.1};
\draw[densely dotted,black!36] (82.312,181.440) -- (137.000,181.440);
\node[anchor=west,inner sep=0pt,font=\rmfamily\fontsize{7}{7}\selectfont] at (138.000,181.440) {EF380031.1};
\draw[densely dotted,black!36] (81.427,184.680) -- (137.000,184.680);
\node[anchor=west,inner sep=0pt,font=\rmfamily\fontsize{7}{7}\selectfont] at (138.000,184.680) {EF380032.1};
\draw[densely dotted,black!36] (102.157,187.920) -- (137.000,187.920);
\node[anchor=west,inner sep=0pt,font=\rmfamily\fontsize{7}{7}\selectfont] at (138.000,187.920) {EF380011.1};
\draw[densely dotted,black!36] (107.649,191.160) -- (137.000,191.160);
\node[anchor=west,inner sep=0pt,font=\rmfamily\fontsize{7}{7}\selectfont] at (138.000,191.160) {EF380012.1};
\draw[densely dotted,black!36] (70.459,194.400) -- (137.000,194.400);
\node[anchor=west,inner sep=0pt,font=\rmfamily\fontsize{7}{7}\selectfont] at (138.000,194.400) {NC\_001422.1};
\draw[densely dotted,black!36] (71.250,197.640) -- (137.000,197.640);
\node[anchor=west,inner sep=0pt,font=\rmfamily\fontsize{7}{7}\selectfont] at (138.000,197.640) {CP004084.1};
\draw[densely dotted,black!36] (63.798,200.880) -- (137.000,200.880);
\node[anchor=west,inner sep=0pt,font=\rmfamily\fontsize{7}{7}\selectfont] at (138.000,200.880) {HM753660.1};
\draw[densely dotted,black!36] (63.798,204.120) -- (137.000,204.120);
\node[anchor=west,inner sep=0pt,font=\rmfamily\fontsize{7}{7}\selectfont] at (138.000,204.120) {HM753661.1};
\draw[line width=.6pt] (62.880,210) -- ++(9.045,0) node[midway,below=2pt,font=\scriptsize] {10 éditions};
\end{tikzpicture}
}
\caption{Les deux phylogrammes portent les identifiants des séquences et une barre horizontale de dix éditions pour chaque arbre.}\label{fig:arbres}\sourcefig{report/phylograms.py}
\end{figure}
\synthese{L'écart absolu maximal vaut 56,87 éditions pour UPGMA et 63,86 pour NJ.}{L'erreur maximale et l'ajustement moyen de la figure~\ref{fig:phage} ordonnent donc les arbres différemment.}

\chapitre{Structure des groupes dans les arbres de surveillance}
Les distances entre 64 séquences archivées de SARS-CoV-2 permettent d'examiner si les groupes étiquetés sont séparés par une arête. Dans un arbre enraciné, un groupe est retrouvé si les descendants d'une arête positive coïncident avec ses membres. Pour NJ, non enraciné, les deux côtés de l'arête sont examinés. Sur les 1\,905\,372 caractères lus, 65\,588 sont ambigus. La distance compte les différences entre caractères déposés, y compris aux positions ambiguës.
\chart{19_phylogenie_surveillance.png}{Les distances globales entre paires et les groupes isolés par une arête sont comparés aux partitions de l'arbre de surveillance sur les identifiants communs.}{fig:surveillance}{scripts/phylo_groups.py}
Les 2\,016 scores de NvPD correspondent à ceux d'edlib et de WFA2. Sur M1 Pro, six processus NvPD monofils qui se partagent les paires terminent en 268,83\,s. La commande qui attribue huit fils à chaque paire atteint la limite de 1\,200\,s. À la figure~\ref{fig:surveillance}, la médiane est de 112 éditions pour les 252 paires de même lignage et de 402 pour les 1\,701 paires de lignages différents. Ces statistiques excluent les 63 paires qui contiennent la référence. Sur 52 feuilles communes, l'arbre public UShER~\cite{usher} donne des longueurs de chemin en événements portés par l'arbre. Leur corrélation descriptive avec les distances d'édition vaut \(r=0,124\) sur 1\,326 paires qui partagent des feuilles.
\synthese{Sur la cohorte entière, UPGMA isole deux des sept groupes et NJ en isole trois.}{Sur les 52 feuilles partagées, les arbres recalculés retrouvent respectivement six et huit des 27 partitions internes positives de l'arbre public.}

\chapitre{Régimes de calcul selon la charge}
Le choix de la décomposition dépend du travail par paire et du nombre de paires disponibles. Pour NvPD synchronisé à chaque ligne, sur une paire de \(m\) lignes et \(n\) colonnes réparties entre \(p\) unités, un modèle de durée sépare le calcul des synchronisations,
\[
 T_{\rm dense}\simeq \frac{mn}{p\,v_p}+m\,\tau_p,
 \qquad W_{\rm bande}=O(\max(m,n)(d+1)).
\]
Le débit local \(v_p\) comprend les effets de la mémoire et \(\tau_p\) désigne le coût moyen d'une synchronisation. Le premier terme suppose des colonnes équilibrées. La seconde expression borne le travail d'une bande doublée jusqu'à couvrir la distance \(d\). Pour des paires indépendantes, la somme des travaux individuels remplace \(mn\) et une réduction termine le lot.
\chart{17_synthese.png}{La comparaison associe le coût des méthodes sur une paire à la variation du débit quand les deux nœuds se partagent un lot de taille fixe.}{fig:bilan}{notebooks/profil_nvpd.ipynb}
Pour les paires d'un million de bases de la figure~\ref{fig:bilan}, le rapport des durées NvPD sur GH200 et WFA2 sur CPU passe de 8,74 à 0,48 entre les taux de 2 et 10\,\%. À taille de lot fixe, chaque fil MPI traite toute sa paire sans collectif par ligne.
\synthese{À deux nœuds, le débit atteint 17\,586 paires de phage par seconde et 3\,820 paires synthétiques par seconde.}{Une faible distance réduit le travail du front. Sur un lot, la longueur des paires détermine aussi le débit.}

\clearpage
\phantomsection\addcontentsline{toc}{section}{Références}
\begin{thebibliography}{16}
\bibitem{sadiq} M. U. Sadiq et coll. \href{https://doi.org/10.1007/s10586-019-02962-w}{\emph{NvPD: novel parallel edit distance algorithm, correctness, and performance evaluation.}} Cluster Computing \textbf{23}, 879--894 (2020).
\bibitem{edlib} M. Šošić et M. Šikić. \href{https://doi.org/10.1093/bioinformatics/btw753}{\emph{Edlib: a C/C++ library for fast, exact sequence alignment using edit distance.}} Bioinformatics \textbf{33}, 1394--1395 (2017).
\bibitem{wfa} S. Marco-Sola et coll. \href{https://doi.org/10.1093/bioinformatics/btaa777}{\emph{Fast gap-affine pairwise alignment using the wavefront algorithm.}} Bioinformatics \textbf{37}, 456--463 (2021).
\bibitem{wfagpu} Q. Aguado-Puig et coll. \href{https://doi.org/10.1109/ACCESS.2022.3182714}{\emph{Accelerating Edit-Distance Sequence Alignment on GPU Using the Wavefront Algorithm.}} IEEE Access \textbf{10}, 63782--63796 (2022).
\bibitem{amdahl} G. M. Amdahl. \href{https://doi.org/10.1145/1465482.1465560}{\emph{Validity of the Single Processor Approach to Achieving Large Scale Computing Capabilities.}} AFIPS, 483--485 (1967).
\bibitem{gustafson} J. L. Gustafson. \href{https://doi.org/10.1145/42411.42415}{\emph{Reevaluating Amdahl's Law.}} Communications of the ACM \textbf{31}, 532--533 (1988).
\bibitem{nj} N. Saitou et M. Nei. \href{https://doi.org/10.1093/oxfordjournals.molbev.a040454}{\emph{The neighbor-joining method: a new method for reconstructing phylogenetic trees.}} Molecular Biology and Evolution \textbf{4}, 406--425 (1987).
\bibitem{ukkonen} E. Ukkonen. \href{https://doi.org/10.1016/S0019-9958(85)80046-2}{\emph{Algorithms for approximate string matching.}} Information and Control \textbf{64}, 100--118 (1985).
\bibitem{myers} G. Myers. \href{https://doi.org/10.1145/316542.316550}{\emph{A fast bit-vector algorithm for approximate string matching based on dynamic programming.}} Journal of the ACM \textbf{46}, 395--415 (1999).
\bibitem{quicked} M. Doblas et coll. \href{https://doi.org/10.1093/bioinformatics/btaf112}{\emph{QuickEd: high-performance exact sequence alignment based on bound-and-align.}} Bioinformatics \textbf{41}, btaf112 (2025).
\bibitem{astar} R. Groot Koerkamp. \href{https://doi.org/10.4230/LIPIcs.WABI.2024.17}{\emph{A*PA2: Up to 19× Faster Exact Global Alignment.}} WABI, LIPIcs \textbf{312}, article 17 (2024).
\bibitem{sassy} R. Beeloo et R. Groot Koerkamp. \href{https://doi.org/10.1093/bioinformatics/btag244}{\emph{Sassy: fuzzy searching DNA sequences using SIMD.}} Bioinformatics \textbf{42}, btag244 (2026).
\bibitem{kokkos} H. C. Edwards, C. R. Trott et D. Sunderland. \href{https://doi.org/10.1016/j.jpdc.2014.07.003}{\emph{Kokkos: Enabling manycore performance portability through polymorphic memory access patterns.}} Journal of Parallel and Distributed Computing \textbf{74}, 3202--3216 (2014).
\bibitem{t2t} S. Nurk et coll. \href{https://doi.org/10.1126/science.abj6987}{\emph{The complete sequence of a human genome.}} Science \textbf{376}, 44--53 (2022).
\bibitem{chmreleases} T2T Consortium. \href{https://github.com/marbl/CHM13/blob/master/Earlier_assembly_releases_and_associated_data.md}{\emph{Earlier assembly releases and associated data.}} Versions CHM13v1.0 et CHM13v2.0.
\bibitem{usher} Y. Turakhia et coll. \href{https://doi.org/10.1038/s41588-021-00862-7}{\emph{Ultrafast Sample placement on Existing tRees (UShER) enables real-time phylogenetics for the SARS-CoV-2 pandemic.}} Nature Genetics \textbf{53}, 809--816 (2021).
\bibitem{hirschberg} D. S. Hirschberg. \href{https://doi.org/10.1145/360825.360861}{\emph{A linear space algorithm for computing maximal common subsequences.}} Communications of the ACM \textbf{18}, 341--343 (1975).
\bibitem{upgma} R. R. Sokal et C. D. Michener. \href{https://www.biodiversitylibrary.org/part/385476}{\emph{A statistical method for evaluating systematic relationships.}} University of Kansas Science Bulletin \textbf{38}, 1409--1438 (1958).
\end{thebibliography}

\end{document}
